% concatenated from Bounded_Ratios_for_Lorentzian_Matrices.tex

\documentclass{amsart}

%\usepackage{etex}
\usepackage{amsmath,amssymb,amsthm,amsfonts,mathrsfs,mathtools}
\usepackage[frame,cmtip,arrow,matrix,line,graph,curve]{xy}
%\usepackage{graphpap,color,paralist,pstricks}
\usepackage[mathscr]{eucal}
\usepackage{mathabx}
\usepackage[colorlinks,backref=page,citecolor=blue]{hyperref}
\usepackage{tikz}
\usepackage{epic,eepic}
\usepackage{yfonts}
\usepackage{enumerate}
\usepackage{comment}

\newcommand{\pentagon}{\tikz[scale=0.1,baseline=-0.5ex]
  \draw (90:1) -- (162:1) -- (234:1) -- (306:1) -- (18:1) -- cycle;}


\renewcommand{\familydefault}{ppl}
\setlength{\oddsidemargin}{0in}
\setlength{\evensidemargin}{0in}
\setlength{\marginparwidth}{0in}
\setlength{\marginparsep}{0in}
\setlength{\marginparpush}{0in}
\setlength{\topmargin}{0.3in}
\setlength{\headsep}{14pt}
\setlength{\footskip}{.3in}
\setlength{\textheight}{8.0in}
\setlength{\textwidth}{5.8in}
\setlength{\parskip}{4pt}
\linespread{1.2}

\theoremstyle{definition}
\newtheorem{mainthm}{Theorem}
\newtheorem{mainconjecture}{Conjecture}
\newtheorem{theorem}{Theorem}[section]
\newtheorem{definition}[theorem]{Definition}
\newtheorem{conjecture}[theorem]{Conjecture}
\newtheorem{lemma}[theorem]{Lemma}
\newtheorem{proposition}[theorem]{Proposition}
\newtheorem{corollary}[theorem]{Corollary}
\newtheorem{problem}[theorem]{Problem}
\newtheorem{question}[theorem]{Question}
\newtheorem*{theorem*}{Theorem}
\newtheorem{exercise}{Exercise}
\theoremstyle{remark}
\newtheorem{remark}[theorem]{Remark}
\newtheorem{example}[theorem]{Example}

\def\PP{\mathbb{P}}
\def\RR{\mathbb{R}}
\def\CC{\mathbb{C}}
\def\ZZ{\mathbb{Z}}
\def\QQ{\mathbb{Q}}
\def\T{\mathbb{T}}
\def\L{\mathbb{L}}
\def\BR{\mathrm{BR}}
\def\Sym{\mathrm{Sym}}
\def\Cut{\mathrm{Cut}}


\DeclareMathOperator{\Vol}{Vol}


\begin{document}

\title{Bounded ratios for Lorentzian matrices}

%\author{Daoji Huang, June Huh, Daniel Soskin, and Botong Wang}

\author{Daoji Huang}
\address{\rm Daoji Huang, University of Massachusetts Amherst}
\email{daojihuang@umass.edu}

\author{June Huh}
\address{\rm June Huh, Princeton University and Korea Institute for Advanced Study}
\email{huh@princeton.edu}

\author{Daniel Soskin}
\address{\rm Daniel Soskin, University of California, Los Angeles}
\email{dsoskin@math.ucla.edu}

\author{Botong Wang}
\address{\rm Botong Wang, University of Wisconsin–Madison}
\email{wang@math.wisc.edu}

\begin{abstract}
We study multiplicative inequalities among entries of Lorentzian matrices, referred to as \emph{bounded ratios}.  
These inequalities can be viewed as generalizations of  the classical Alexandrov--Fenchel inequalities for mixed volumes.
Our main structural result identifies the cone of all bounded ratios on Lorentzian matrices with the dual of the \emph{cut cone}, a finitely generated integral polyhedral cone extensively studied in metric geometry and graph theory. 
We examine in detail the \emph{pentagonal ratio}, which first appears for Lorentzian matrices of size at least five.
For Lorentzian matrices of size three, we determine the optimal bounding constants across the entire cone of bounded ratios, obtaining an explicit entropy-like formula. We conjecture that any normalized bounded ratio is, in fact, bounded by $2$.
\end{abstract}
\maketitle

\begin{comment}


-- Introduce the notion of Lorentzian matrices and indicate significance.

-- Define the notion of bounded ratios for Lorentzian matrices.

-- Recall earlier work on positive definite matrices.

-- Recall Brunn-Minkowski theory, and in particular Alexandrov-Fenchel inequality for mixed volumes of convex bodies. These are bounded ratios that appear when n=2.

-- Introduce reverse Khovanskii-Teissier for n=3, and the new ratio for n=5.

-- Theorem A: A complete description of bounded ratios for n=5, including constants.


-- For example, the mixed volumes of five convex bodies (plus auxiliary convex bodies) satisfy that inequality for n=5.

-- We show that the set of bounded ratios is an integral polyhedral cone. In particular, there are finitely many extremal bounded ratios for n-by-n Lorentzian matrices. Up to $S_n$-symmetry, the number of extremal bounded ratios are 1,1,2,4,36, 2169,...



-------------------------------------------------------------------

-- Define the cut-cone as the space of $\ell_1$-semimetrics.
Read the book on the cut-cone and summarize the main results and their significance. 

-- We reveal the main result.

Theorem B: The cone of bounded ratios for the n-by-n Lorentzian matrices is the dual of the cut-cone.

-- For example, hypermetrics always give bounded ratios. Say that $n=5$ example above is hypermetric. Maybe two new ones for $n=6$.


-------------------------------------------------------------------





-- We now summarize the main technique in the broader context.

-- (1) Briefly summarize Lorentzian polynomials [Branden-Huh].

-- (2) Briefly summarize matroid representations over tracts [Baker-Bowler].

-- (3) Briefly summarize polymatroid representations over tracts. This is a manuscript under preparation [Baker-Huh-Kummer-Lorscheid], attached. It's under construction, but it should be readable. Please don't share.

Our techniques are motivated by the general observation that
\[
\log \mathbb{T}_0 \subseteq \log \mathbb{L} \subseteq \log \mathbb{T}_2 \subseteq \text{$\delta$-hyperbolic spaces}
\]

-- (5) (Reading required) Briefly summarize Gromov's work on hyperbolic groups and spaces. State the tree approximation theorem.


We now have the language to outline our main argument concisely. Sketch the argument. 


-------------------------------------------------------------------

The last part of the introduction will be on conjectures about bounding constants.

State the end result of the analysis for $n=3$, including a picture of a triangle with inscribed circle. Outside the circle, the function is given by the following entropy-like expression
\[
xxx
\]

-- Conjecture: The normalized bounded ratios are bounded by $2$.

Discussion of rank $2$ Lorentzian matrices.


-------------------------------------------------------------------
\newpage
\end{comment}
\section{introduction}

A \emph{Lorentzian matrix} is a symmetric matrix with nonnegative real entries that has at most one positive eigenvalue.
We denote by $\L_n$ the set of $n\times n$ Lorentzian matrices.
These matrices appear prominently in space-time geometry and convex geometry. In this work, we study \emph{multiplicative inequalities} for Lorentzian matrices---inequalities of the form
\[
\prod_{i \le j} p_{ij}^{\alpha_{ij}} \le c \prod_{i \le j} p_{ij}^{\beta_{ij}} \ \ \text{for all $(p_{ij}) \in \L_n$,}
\]
where $\alpha_{ij}$ and $\beta_{ij}$ are nonnegative integers, and $c$ is a positive constant.

For a collection of convex bodies $K=(K_1, \dots, K_d)$ in $\RR^d$, the \emph{mixed volume} of $K$ is defined to be the normalized coefficient of $t_1\cdots t_d$ in the associated volume polynomial:
\[
V(K_1, \dots, K_d)\coloneq \frac{1}{d!}\partial_1\cdots \partial_d \Vol_n(t_1K_1+\dots t_dK_d).
\]
The celebrated \emph{Alexandrov--Fenchel inequality} states that,
for any collection of convex bodies 
\[
 P_1,\ldots,P_n \ \ \text{and} \ \ K_1,\ldots,K_{d-2} \ \ \text{in $\RR^d$,}
\]
the symmetric $n \times n$ matrix of $d$-dimensional mixed volumes $(p_{ij})$ given by
\[
p_{ij}=V(P_i, P_j, K_1, \dots, K_{d-2}) %\ \ \text{for all $i$ and $j$}
\]
is a Lorentzian matrix \cite[Section 7.3]{Schneider}. In particular, any of its $2 \times 2$ principal minor is nonpositive, that is, 
\[
p_{ii}p_{jj} \le p_{ij}^2 \ \ \text{for all $i$ and $j$.}
\]
%\[
% V(P_i, P_i, K_1, \dots, K_{d-2})V(P_j, P_j, K_1, \dots, K_{d-2}) \le V(P_i, P_j, K_1, \dots, K_{d-2})^2. %\ \ \text{for all $i$ and $j$.}
%\]


The Alexandrov--Fenchel inequality above is one of the first examples of bounded ratios on Lorentzian matrices, which we now define. Let $\Sym_n(\RR_{> 0})$ denote the set of $n \times n$ symmetric matrices with positive entries, and let $X \subseteq \Sym_n(\RR_{> 0})$ be a nonempty subset.
Given a real vector $\alpha=(\alpha_{ij})_{1 \le  i \le j \le n}$, we write
\[
\alpha=\sum_{i \le j} \alpha_{ij} \mathbf{e}_{ij},
\]
where $\{\mathbf{e}_{ij}\}_{i \le j}$ is the standard basis of the space of all such vectors. 

\begin{definition}
The set of \emph{bounded ratios} on $X$, denoted $\BR(X)$, is the set of all real vectors $\alpha=(\alpha_{ij})_{1 \le  i \le j \le n}$
that satisfy the following condition:
\[
\text{There is a positive constant $c$ such that $\prod_{1\le i\le j\le n}p_{ij}^{\alpha_{ij}}\le c$ for all matrices $(p_{ij})$ in $X$.}
\]
We define the \emph{optimal bounding constant} of a bounded ratio $\alpha$ on $X$, denoted $f(\alpha)=f_X(\alpha)$,  to be the infimum of all possible constants $c$ satisfying the displayed inequality for all $(p_{ij})$ in $X$. 
\end{definition}

It is straightforward to check that $\BR(X)$ is a convex cone in $\RR^{\binom{n+1}{2}}$ and $f(\alpha)$ is a log-convex homogeneous function of degree $1$ on the cone of bounded ratios. 
When $X$ is the set of $n\times n$ positive semidefinite matrices with positive entries $\operatorname{PSD}^+_n$, %$\operatorname{PSD}^+_n$, it is a consequence of the work of 
Yu shows in \cite{Yu} that the cone of bounded ratios on $X$  is a finitely generated integral polyhedral cone with $\binom{n}{2}$ extremal rays corresponding to the $2 \times 2$ principal minors:
\[
\BR(\operatorname{PSD}^+_n)=\operatorname{Cone}(-\mathbf{e}_{ii}-\mathbf{e}_{jj}+2\mathbf{e}_{ij})_{1 \le i<j \le n}.
\]
%$\{r_{kl}\}_{1\le k<l\le n}$, $r_{kl}=(\alpha_{ij})_{1\le %i\le j\le n}$ where $\alpha_{kk}=\alpha_{ll}=-1$, %$\alpha_{kl}=2$, and $\alpha_{ij}=0$ otherwise.
In contrast, bounded ratios for Lorentzian matrices, as we will see, exhibit more intricate behaviors. 
For related notions of bounded ratios in the study of cluster algebras and total positivity, see \cite{Fallat-Johnson,Fallat-Gekhtman-Johnson,Gekhtman-Greenberg-Soskin,Gekhtman-Soskin}.

\begin{definition}
A bounded ratio $\alpha$ is \emph{integral} if all the entries $\alpha_{ij}$ are integers. 
An integral bounded ratio on $X$ is said to be \emph{primitive} if it cannot be expressed as the sum of two nonzero integral bounded ratios on $X$.
\end{definition}

Let $\L_n^+$ be the set of $n\times n$ Lorentzian matrices with only positive entries. Since any Lorentzian matrix is a limit of Lorentzian matrices with positive entries \cite[Section 2]{Branden-Huh}, the multiplicative inequalities for $\L_n$ are obtained by clearing the denominator in $p^\alpha \le c$ for a bounded ratio $\alpha$ on $\L_n^+$.
Our main result, Theorem~\ref{mth:BR} below, implies that $\BR(\L_n^+)$ is a finitely generated integral polyhedral cone. 
%We write $f_\alpha$ for optimal constants of $\L_n^+$.

\begin{example}
The bounded ratios on $\L_2^+$ are exactly the nonnegative multiples of 
\[
\alpha^{12}\coloneq \mathbf{e}_{11}+\mathbf{e}_{22}-2\mathbf{e}_{12}.%(\alpha_{11},\alpha_{12},\alpha_{22})=(1,-2,1).
\]
In other words, $\BR(\L_2^+)=\operatorname{Cone}(\alpha^{12})$. 
The optimal bounding constants are determined by the condition $f(\alpha^{12})=1$. We refer to the bounded ratio $\alpha^{12}$ as the \emph{Alexandrov--Fenchel type}.
\end{example}

\begin{example}
Theorem~\ref{mth:BR} shows that 
there are precisely three primitive bounded ratios on $\L_3^+$, corresponding to the inequalities
\[
\frac{p_{23}p_{11}}{p_{12}p_{13}} \le 2, \quad
\frac{p_{13}p_{22}}{p_{12}p_{23}} \le 2, \quad
\frac{p_{12}p_{33}}{p_{13}p_{23}} \le 2 \quad \text{for all $(p_{ij}) \in \L_3^+$.}
\]
The cone of bounded ratios $\BR(\L_3^+)$ is the three-dimensional simplicial cone in $\RR^6$ generated by
\[
\alpha^{23|1}\coloneq \mathbf{e}_{23}+\mathbf{e}_{11}-\mathbf{e}_{12}-\mathbf{e}_{13}, \quad
\alpha^{13|2}\coloneq\mathbf{e}_{13}+\mathbf{e}_{22}-\mathbf{e}_{12}-\mathbf{e}_{23},  \quad
\alpha^{12|3}\coloneq\mathbf{e}_{12}+\mathbf{e}_{33}-\mathbf{e}_{13}-\mathbf{e}_{23}.
\]
We refer to the bounded ratios $\alpha^{ij|k}$ as the \emph{triangular ratios}. The Alexandrov--Fenchel type ratios can be expressed as the pairwise sums
\[
\alpha^{12}=\alpha^{23|1}+\alpha^{13|2},
 \quad
\alpha^{13}=\alpha^{23|1}+\alpha^{12|3},
\quad
\alpha^{23}=\alpha^{13|2}+\alpha^{12|3},
\]
which are bounded but not primitive on $\L_3^+$.
To see the validity of the displayed inequalities on $\L_3^+$, we use the implication
\[
\text{$
\left(
\begin{matrix}
	p_{11} & p_{12} & p_{13} \\
	p_{12} & p_{22} & p_{23} \\
	p_{13} & p_{23} & p_{33}
\end{matrix}\right)$ is Lorentzian}
\Longrightarrow
\text{$
\left(
\begin{matrix}
	0 & p_{12} & p_{13} \\
	p_{12} & 0 & p_{23} \\
	p_{13} & p_{23} & p_{33}
\end{matrix}\right)$ is Lorentzian.}
\]
The nonnegativity of the determinant of the latter matrix gives the optimal bounding constant $f(\alpha^{12|3})=2$, which is witnessed by the Lorentzian matrices
\[
\left(
\begin{matrix}
	\epsilon & 1 & 1 \\
	1 & \epsilon & 1 \\
	1 & 1 & 1
\end{matrix}\right) \in \L_3^+ \ \ \text{for sufficiently small $\epsilon >0$.}
\]
Note that $0= \log f(\alpha^{12|3}+\alpha^{13|2})\le \log f(\alpha^{12|3})+\log f(\alpha^{13|2})=2 \log 2$.
In Theorem~\ref{mth:entropy}, we explicitly describe the function $f(\alpha)$ on $\BR(\L_3^+)$.


The inequalities for the primitive bounded ratios $\alpha^{ij|k}$
 appear in the work of Andr\'e Weil on the proof of the Riemann hypothesis for algebraic curves over finite fields, where they are referred to as the \emph{Castelnuovo--Severi inequality} \cite[Theorem 1.5]{Milne}.
This is the first special case of the \emph{reverse Khovanskii--Teissier inequality} \cite[Theorem 5.7]{Lehmann-Xiao}. The inequalities for $\alpha^{ij|k}$ also appear
in early works on Brunn--Minkowski theory; see \cite[Page 396]{Frobenius1915} and \cite[Section~51]{Bonnesen-Fenchel}.\footnote{We thank Ramon van Handel for pointing out these references.}
\end{example}


It turns out that there are no essentially new bounded ratios on $\L_4^+$:
\[
\BR(\L^+_4)=\operatorname{Cone}(\alpha^{ij|k})_{1 \le i,j,k \le 4}, \ \ \text{where $\alpha^{ij|k} \coloneq \mathbf{e}_{ij}+\mathbf{e}_{kk}-\mathbf{e}_{ik}-\mathbf{e}_{jk}$.}
\]
Exactly one new type of bounded ratios appear for $\L_5^+$, which we refer to as the \emph{pentagonal ratios}:
\[
\alpha^{ijk|lm}\coloneq \mathbf{e}_{ij}+\mathbf{e}_{ik}+\mathbf{e}_{jk}+\mathbf{e}_{ll}+\mathbf{e}_{lm}+\mathbf{e}_{mm}-\mathbf{e}_{il}-\mathbf{e}_{jl}-\mathbf{e}_{kl}-\mathbf{e}_{im}-\mathbf{e}_{jm}-\mathbf{e}_{km}.
\]



\renewcommand{\themainthm}{A}
 
\begin{mainthm}
\label{mth:pentagonal}
The \emph{pentagonal inequality} holds for any $(p_{ij}) \in \L_n$:
\[
p_{ij}p_{ik}p_{jk}p_{ll}p_{lm}p_{mm} \le 4 p_{il}p_{jl}p_{kl}p_{im}p_{jm}p_{km} \ \ \text{for any $i,j,k,l,m$.}
\]
The constant $4$ is optimal when all five indices are distinct, and
    \[
\BR(\L_5^+)=\operatorname{Cone}(\alpha^{ijk|lm})_{1 \le i,j,k,l,m \le 5}.
    \]
\end{mainthm} 

Note that the triangular ratios appear as degenerate cases of the pentagonal ratios:
\[
\alpha^{ijk|kk} = \alpha^{ij|k} \ \ \text{for any $i,j,k$.}
\]
The Alexandrov--Fenchel type ratio $\alpha^{ij}$ appears as a degenerate triangular ratio $\alpha^{ii|j}$, which is bounded but not extremal in $\BR(\L_5^+)$.
The cone of bounded ratios $\BR(\L_5^+)$ has forty extremal rays, corresponding to the ten nondegenerate pentagonal ratios and the thirty nondegenerate triangular ratios.
We have a factorization of the pentagonal ratio  into a product of three  triangular ratios and one reciprocal:
\[
\frac{p_{12}p_{13}p_{23}p_{44}p_{45}p_{55}}{p_{14}p_{15}p_{24}p_{25}p_{34}p_{35}} = \frac{p_{12}p_{44}}{p_{14}p_{24}}\cdot\frac{p_{13}p_{55}}{p_{15}p_{35}}\cdot\frac{p_{22}p_{45}}{p_{24}p_{25}}\cdot\left(\frac{p_{22}p_{34}}{p_{23}p_{24}}\right)^{-1}.
\]
 Even though the reciprocal on the right-hand side is unbounded  on its own, the pentagonal ratio on the left-hand side is bounded on $\L_n^+$. Factorizations of this type are not unique. 


In general, we show that $\BR(\L_n^+)$  is a finitely generated polyhedral cone whose extremal rays are generated by the primitive bounded ratios. Up to $S_n$-symmetry, the numbers of primitive bounded ratios for $n=3,4,5,6,7,8$ are $1$, $1$, $2$, $4$, $36$, and $2169$, respectively.
The two new types of primitive bounded ratios for $\L_6^+$ correspond to the inequalities of the form
\[
\frac{p_{11}p_{12}p_{12}p_{13}p_{13}p_{23}p_{44}p_{45}p_{46}p_{55}p_{56}p_{66}}{p_{14}p_{14}p_{15}p_{15}p_{16}p_{16}p_{24}p_{25}p_{26}p_{34}p_{35}p_{36}} \le c, \quad
\frac{p_{11}p_{11}p_{11}p_{12}p_{12}p_{22}p_{34}p_{35}p_{36}p_{45}p_{46}p_{56}}{p_{13}p_{13}p_{14}p_{14}p_{15}p_{15}p_{16}p_{16}p_{23}p_{24}p_{25}p_{26}} \le c,
\]
although we do not know the optimal bounding constants for these inequalities; see Conjecture~\ref{conj:intro-at-most-2}.  
To prove the general statement, we relate $\BR(\L_n^+)$ to the \emph{cut cone}, a finitely generated integral polyhedral cone extensively studied in metric geometry and graph theory. 
For a comprehensive introduction to the geometry and combinatorics of the cut cone, we refer to the book \cite{Deza-Laurent}.

\begin{definition}
For any subset $S \subseteq [n]$, define the \emph{cut vector} $\delta(S)=(\delta(S)_{ij})_{1 \le i < j \le n}$ by
\[
\delta(S)_{ij} = 
\begin{cases}
1, & \text{if $S$ contains exactly one of $i$ and $j$}, \\
0, & \text{if otherwise.}
\end{cases}
\]
The \emph{cut cone} $\Cut_n$ is the polyhedral cone in $\RR^{\binom{n}{2}}$ generated by all the cut vectors $\delta(S)$ for $S \subseteq [n]$.
\end{definition}

The extremal rays of $\Cut_n$ are generated by the cut vectors of nonempty proper subsets of $[n]$.
The cut cone can be identified with the space of semimetrics on $n$ points that are isometrically embeddable in $\ell_1$-space \cite[Section 1.1]{Deza-Laurent}. 

We consider the \emph{scaling action} of 
$(\RR_{>0})^n$ on $\Sym_n(\RR_{> 0})$ defined by the formula
%\begin{equation}
%\label{eq:torus-action}
\[
	(c_1,\dots, c_n)\cdot (p_{ij}) \coloneqq (c_ic_jp_{ij}).
    \]
%\end{equation}
If $X\subseteq \Sym_n(\RR_{>0})$ is stable under the scaling action, then the orbit space of $X$ can be identified with the set $\underline{X}$ of symmetric matrices in $X$ with diagonal entries $1$.
%\[
%\underline{X}\coloneqq \Big\{(p_{ij}) \in X, \ \  \text{$p_{ii}=1$ for all $i$} \ \Big\}.
%\]
We say that such matrices in $\underline{X} \subseteq X$ are \emph{normalized}, and define the cone of  \emph{reduced bounded ratios} on $X$ by 
\[
\underline{\BR}(X)
%\BR'(\underline{X})
\coloneqq
\left\{ 
(\alpha_{ij})_{1 \le i<j \le n} \;\middle|\; \text{there is $c>0$ such that  $\prod_{1\le i< j\le n}p_{ij}^{\alpha_{ij}}\le c$ for all $(p_{ij}) \in \underline{X}$}   \right\}.
\]
As before, the optimal bounding constants $\underline{f}(\alpha)$ for reduced bounded ratios are defined to be the infimum of all possible constants $c$ satisfying the displayed inequality for all $(p_{ij})$ in $\underline{X}$.
It is straightforward to check that $\underline{\BR}(X)$ is a convex cone in $\RR^{\binom{n}{2}}$ and $\underline{f}(\alpha)$ is a log-convex homogeneous function of degree $1$ on the cone of reduced bounded ratios. 


%Constraints on bounded ratios are revealed by the scaling action: 
If $\alpha$ is a bounded ratio on $X$, then the corresponding monomial $p^\alpha$
must be invariant under the scaling action of $(\RR_{>0})^n$. 
In other words, if $\alpha\in\BR(X)$, then 
\[
%\label{eq:homogenize-alpha}
2\alpha_{ii}=-\sum_{j\neq i}\alpha_{ij} \ \ \text{for all $i$}.
\]
Therefore, the projection $\pi: \RR^{\binom{n+1}{2}} \to \RR^{\binom{n}{2}}$ that omits the diagonal entries restricts to a bijection between $\BR(X)$ and $\underline{\BR}(X)$. Moreover,  we have $f(\alpha)=\underline{f}(\pi(\alpha))$ for any bounded ratio $\alpha$.



\renewcommand{\themainthm}{B}
\begin{mainthm}%[Theorem~\ref{thm:BR}]
\label{mth:BR}
	The cone of reduced bounded ratios on $\L_n^+$ is %$\underline{\BR}(\L_n^+)$
     dual of the cut cone:
	\[
    \underline{\BR}(\L_n^+)=\left\{(\alpha_{ij})_{1 \le i < j\le n} \;\middle|\; \sum_{1 \le i<j \le n} \alpha_{ij} \beta_{ij} \le 0 \ \ \text{for all $(\beta_{ij})_{1 \le i<j \le n} \in \Cut_n$} \right\} \subseteq \RR^{\binom{n}{2}}.
    \]
In particular, the cone of bounded ratios on $\L_n^+$ is a  finitely generated integral polyhedral cone.
\end{mainthm}

Thus, the primitive bounded ratios on $\L_n^+$ is in bijection with the facets of $\Cut_n$. Karp and Papadimitriou showed in  \cite{Karp-Papadimitriou}  that  there is no computationally tractable description of all facets of the cut cone  unless $\mathrm{NP}=\mathrm{coNP}$.
Consequently, obtaining a complete and explicit description of the primitive bounded ratios on $\L_n^+$ is likely to be a challenging problem.


Nevertheless, several interesting classes of facets of $\Cut_n$ are known, the first among which is given by  the \emph{hypermetric inequalities} \cite[Chapter 28]{Deza-Laurent}: For any vector $h=(h_i)_{1 \le i \le n}$ satisfying $\sum_{i=1}^n h_i=1$, we have
\[
\sum_{1 \le i < j \le n} h_ih_j \beta_{ij} \le 0 \ \ \text{for all $(\beta_{ij})_{1 \le i<j \le n}$ in $\Cut_n$.}
\]
Thus, we may deduce the following statement from Theorem~\ref{mth:BR}.

\begin{corollary}
If  $\sum_{i=1}^n h_i=1$, the vector $(h_ih_j)_{1 \le i<j\le n}$ is a reduced bounded ratio on $\L_n^+$.
\end{corollary}

Bounded ratios of this type are called \emph{hypermetric}. Determining the optimal bounding constants for hypermetric bounded ratios as a function of $h$ is an interesting and challenging problem.
The hypermetric inequality for $h=\mathbf{e}_1+\mathbf{e}_2-\mathbf{e}_3$ gives the triangular ratio $\alpha^{12|3}$, and the hypermetric inequality for $h=\mathbf{e}_1+\mathbf{e}_2+\mathbf{e}_3-\mathbf{e}_4-\mathbf{e}_5$ gives the pentagonal ratio $\alpha^{123|45}$.
Up to $S_n$-symmetry, two additional hypermetric inequalities are needed to describe all facets of $\Cut_6$ \cite[Remark 15.2.11]{Deza-Laurent}. These hypermetric inequalities correspond to 
\[
h=2\mathbf{e}_1+\mathbf{e}_2+\mathbf{e}_3-\mathbf{e}_4-\mathbf{e}_5-\mathbf{e}_6  \ \ \text{and} \ \ 
h=-2\mathbf{e}_1-\mathbf{e}_2+\mathbf{e}_3+\mathbf{e}_4+\mathbf{e}_5+\mathbf{e}_6.
\]
For $n \ge 7$, the cut cone has facets that do not correspond to any hypermetric inequality; see \cite[Section 30.6]{Deza-Laurent} for a complete list of facets of $\Cut_7$. 
For a comprehensive discussion of the facets of the cut cone, we refer the reader to Deza and Laurent’s monograph \cite[Part V]{Deza-Laurent}.

We begin by outlining the proof techniques for our main theorem and situating them within a broader context. Br\"and\'en and the second author introduced \emph{Lorentzian polynomials} \cite{Branden-Huh}, a class of homogeneous polynomials that unify and extend various notions of log-concavity arising in combinatorics and geometry. A quadratic Lorentzian polynomial is, by definition, a quadratic form whose Hessian matrix is Lorentzian. 
 Meanwhile, Baker and Bowler developed a unifying algebraic framework for \emph{matroids over tracts} \cite{Baker-Bowler}, which simultaneously generalized many variants of matroids. Building on this framework, the authors of  \cite{BHKL-triangular} identified the space of Lorentzian polynomials with given support with the corresponding thin Schubert cell in the Grassmannian over the \emph{triangular hyperfield}, up to homeomorphism. This enabled the authors to use Gromov's theorem on $\delta$-hyperbolic spaces in \cite{BHKL-triangular2} to derive new results on Lorentzian polynomials. Our proof is inspired by this perspective, which we now explain in concrete terms.

 Let $\Sym_n(\RR_{ \ge 0})$ denote the set of $n \times n$ symmetric matrices with nonnegative entries.
 
\begin{definition}
Let $\Delta_n(\T_0)$ be the set of matrices $(p_{ij}) \in \Sym_n(\RR_{ \ge 0})$ such that
\[
\text{the maximum among $p_{ij}p_{kl}, p_{ik}p_{jl}, p_{il}p_{jk}$ is achieved at least twice for any $i,j,k,l \in [n]$.}
\]
For a positive real number $p$, let $\Delta_n(\T_p)$ be the set of matrices $(p_{ij}) \in \Sym_n(\RR_{ \ge 0})$ such that
\[
\text{
$(p_{ij}p_{kl})^{1/p} \le
 (p_{ik}p_{jl})^{1/p}+
(p_{il}p_{jk})^{1/p}$  for any $i,j,k,l \in [n]$.}
\]
For any nonnegative real number $p$, let $\Delta^+_n(\T_p)$ denote the intersection $\Delta_n(\T_p) \cap \Sym_n(\RR_{> 0})$.
\end{definition}



%In the language of \cite{BHKL-triangular}, $\Delta_n(\T_p)$ is the affine cone over the Grassmannian $\mathrm{Gr}(2,n)(\T_p)$ over the triangular hyperfield $\T_p$. For our purposes the explicit definition given above suffices.

Note that $\Delta^+_n(\T_p)$ is invariant under the scaling action of $(\RR_{>0})^n$. As before, we identify the orbit space with the set $\underline{\Delta}^+_n(\T_p)$ of matrices in $\Delta^+_n(\T_p)$ with diagonal entries $1$.
A key observation from \cite{BHKL-triangular} is that 
\[
 \underline{\Delta}_n^+(\T_0)\subseteq  \underline{\L}_n^+\subseteq  \underline{\Delta}_n^+(\T_2),
\]
where $\underline{\L}_n^+$ is the set of $n \times n$  Lorentzian matrices with diagonal entries $1$ and off-diagonal entries positive.
By taking the coordinatewise logarithm and omitting the diagonal coordinates, we have
\[
 \log \underline{\Delta}_n^+(\T_0)\subseteq  \log \underline{\L}_n^+\subseteq  \log \underline{\Delta}_n^+(\T_2) \ \ \text{in} \ \  \RR^{\binom{n}{2}}.
\]
We observe in Section~\ref{sec:hyperbolic} that  $\log \underline{\Delta}_n^+(\T_2)$ parametrizes Gromov's $\delta$-hyperbolic metrics on $n$ points for $\delta= \log 2$, while $ \log \underline{\Delta}_n^+(\T_0)$ parametrizes $0$-hyperbolic metrics on $n$ points, which are precisely the tree metrics. It follows from Gromov's tree approximation theorem \cite[Section 6.1]{Gromov} 
%and \cite[90 Section 2]{GdlH} 
that the two spaces %in $ \log \underline{\Delta}_n^+(\T_0)\subseteq  \log \underline{\L}_n^+\subseteq  \log \underline{\Delta}_n^+(\T_2)$ 
are within finite Hausdorff distance of one another. Consequently, 
\[
\underline{\BR}(\Delta_n^+(\T_0))=\underline{\BR}(\L_n^+)=\underline{\BR}(\Delta_n^+(\T_p)).
\]
Since $\Cut_n$ is the convex hull of the space of tree metrics on $n$ points, we conclude that $\underline{\BR}(\L_n^+)$ and $\Cut_n$ are dual to each other; see Section~\ref{sec:hyperbolic} for further details.



Having understood the cone of bounded ratios for Lorentzian matrices, it is natural  to ask how the optimal bounding constants $f(\alpha)$ behave as a function on this cone. As mentioned before,
\[
\log f(\alpha+\beta) \le \log f(\alpha)+\log f(\beta) \ \ \text{and} \ \ \log f(c \alpha) = c \log f(\alpha).
\]
In Theorem~\ref{mth:entropy} below, we give an explicit description of the  function $f(\alpha)$ when $n=3$.

Recall that the cone of bounded ratios $\BR(\L_3^+)$ has three extremal rays 
$
\alpha^{23|1}, \alpha^{13|2}, \alpha^{12|3}
$
corresponding to the optimal inequalities
\[
\frac{p_{23}p_{11}}{p_{12}p_{13}} \le 2, \quad \frac{p_{13}p_{22}}{p_{12}p_{23}} \le 2,
\quad \frac{p_{12}p_{33}}{p_{13}p_{23}} \le 2.
\]
Since $\log f$ is a homogeneous function of degree $1$, it is enough to determine 
\[
f(a,b,c)\coloneq f(a \cdot \alpha^{23|1}+b \cdot \alpha^{13|2}+ c \cdot \alpha^{12|3}) \ \ \text{when  $a+b+c=1$ and $a,b,c \ge 0$.}
\]


\renewcommand{\themainthm}{C}
\begin{mainthm}
\label{mth:entropy}
Let $a,b,c$ be nonnegative numbers such that $a+b+c=1$.
\begin{enumerate}[(1)]\itemsep 5pt
\item If $a^2+b^2+c^2-2ab-2ac-2bc\le 0$, then $ f(a,b,c)=1$.
\item If $a^2+b^2+c^2-2ab-2ac-2bc\ge 0$, $a \ge b$,  $a \ge c$, then
\[
f(a,b,c)=2\cdot a^a\cdot b^b \cdot c^c \cdot (2a-1)^{2a-1} \cdot (1-2b)^{2b-1}\cdot (1-2c)^{2c-1}.
\]
\end{enumerate}
See Figure~\ref{fig:triangle} for an illustration.
\end{mainthm}

\begin{figure}[htbp]
    \centering
    \includegraphics[width=0.6\linewidth]{equilateral_triangle.png}
    \caption{Optimal bounding constants for an equilateral triangle cross section of $\underline{\BR}(\L_3^+)$. The midpoints of the three edges correspond to the Alexandrov--Fenchel type inequalities. Inside the inscribed circle, the optimal bounding constant is $1$. Outside of the inscribed circle, the optimal bounding  constants are given by an entropy-like function.}
    \label{fig:triangle}
\end{figure}

We conclude with a general conjecture on the optimal bounding constant $f(\alpha)$.

\begin{definition}
We say that a reduced bounded ratio on $\L_n^+$ is \emph{normalized} if the sum of its coordinate is $-1$.
\end{definition}

For example, reduced bounded ratios of the form  $\pi(\alpha^{ij|k})=\mathbf{e}_{ij}-\mathbf{e}_{ik}-\mathbf{e}_{jk}$ are normalized. 
Since the set of Lorentzian matrices is preserved when a diagonal entry is replaced by $0$,
the coordinate sum of any nonzero reduced bounded ratio is negative. %see Lemma~\ref{lem:alpha-negative}. 
Thus, any nonzero reduced bounded ratio is uniquely a positive multiple of a normalized reduced bounded ratio.

%\renewcommand{\themainconjecture}{D}
\begin{conjecture}%[Conjecture~\ref{conj:at-most-2}]
\label{conj:intro-at-most-2}
	The optimal bounding constant for any normalized reduced bounded ratio on $\L_n^+$ is at most 2.
\end{conjecture} 

Using log-convexity of the optimal bounding constants, 
Conjecture~\ref{conj:intro-at-most-2} can be reduced to the case of extremal bounded ratios. Thus,
by Theorem~\ref{mth:pentagonal}, we know that Conjecture \ref{conj:intro-at-most-2} holds for $n\leq 5$.
See Section~\ref{sec:rank2} for a closely related conjecture for rank $2$ Lorentzian matrices.


%Finally, we analyze the locus of Lorentzian matrices of rank at most 2, denoted $\L_n^+|_{\mathrm{rk2}}$, and give a strengthened conjecture about bounded ratios. Lorentzian matrices of rank at most 2 correspond to quadratic forms that can be factored as a product of two linear forms. We show that  $\BR'(\underline{\L_n|_{\mathrm{rk2}}})$ is also the dual of the cut cone. We make the following conjecture.
%\renewcommand{\themainconjecture}{E}
%\begin{mainconjecture}[Conjecture~\ref{conj:sub-free}]
%\label{conj:intro-sub-free}
%	Let $\alpha=(\alpha_{ij})_{1\le i\le j\le n}$ be a primitive bounded ratio of $\BR(\L_n^+|_\mathrm{rk2})=\BR(\L^+_n)$. Let $p_{ij}=\frac{1}{2}(a_ib_j+a_jb_i)$ for all $i,j\in[n]$. Then 
%	\[2^{\sum_{i=1}^n\alpha_{ii}}\Biggl(\prod_{\substack{i\le j\\\alpha_{ij}<0}}p_{ij}^{-\alpha_{ij}}\Biggr)-\Biggl(\prod_{\substack{i\le j\\\alpha_{ij}>0}}p_{ij}^{\alpha_{ij}}\Biggr)\]
%	is a subtraction-free expression in monomials of $a_i$ and $b_i$. 
%\end{mainconjecture}

\subsection*{Acknowledgements}
The authors gratefully acknowledge the Institute for Advanced Study for providing an inspiring research environment and thank Matt Baker, Chayim Lowen, and Ramon van Handel for insightful and stimulating discussions.
Daoji Huang is  supported by the Charles Simonyi Endowment and NSF-DMS2202900.
June Huh is partially supported by the Simons Investigator Grant.
 Botong Wang is partially supported by the NSF grant DMS-1926686.

 \section{The pentagonal inequality}


 
We prove Theorem~\ref{mth:pentagonal}.
The primary difficulty in proving this theorem lies in determining the optimal bounding constant for the pentagonal inequality. The statement that $\BR(\L_5^+)$ is generated by the pentagonal ratios follows from Theorem~\ref{mth:BR} and the fact that the facets of $\Cut_5$ are defined by the pentagonal ratios \cite[Section 30.6]{Deza-Laurent}.

We begin by observing that the constant $4$ in the pentagonal inequality is, if valid, best possible.
Let $t$ be a nonnegative real number, and consider the matrix
\[
M(t)\coloneq
\begin{pmatrix}
0 & 1 & 1 & t & 2 + t \\[1pt]
1 & 0 & 1 & 2 & 2 \\[1pt]
1 & 1 & 0 & 2 + t & t \\[1pt]
t & 2 & 2 + t & 4t & 4 + 4t \\[1pt]
2 + t & 2 & t & 4 + 4t & 4t
\end{pmatrix}.
\]
It is straightforward to check that $M(t)$ has rank $3$ for $t \ge 0$.
Since the leading $3 \times 3$ principal submatrix of $M$ is nondegenerate and Lorentzian, we see that $M(t)$ is Lorentzian for any $t \ge 0$. We have
\[
\lim_{t \to 0} \frac{p_{12}p_{13}p_{23}p_{44}p_{45}p_{55}}{p_{14}p_{15}p_{24}p_{25}p_{34}p_{35}}=\lim_{t \to 0} \frac{(4t)^2(4+4t)}{4t^2(2+t)^2}=4. 
\]
Since any Lorentzian matrix is a limit of Lorentzian matrices with positive entries \cite[Section~2]{Branden-Huh},
%Since the pentagonal ratio continuously extends to Lorentzian matrices with $p_{11}=p_{22}=p_{33}=0$ and the remaining entries positive, 
there is a family in $\L_5^+$
on which the evaluation of the pentagonal ratio limits to $4$.

We prepare the proof of the pentagonal inequality with a few auxiliary lemmas.

\begin{lemma}\label{lem:factors}
Suppose that $A$ is a symmetric matrix of the form
\[
A=\begin{pmatrix}
0 & 1 & 1 & p_{14} & p_{15} \\
1 & 0 & 1 & p_{24} & p_{25} \\
1 & 1 & 0 & p_{34} & p_{35} \\
p_{14} & p_{24} & p_{34} & p_{44} & p_{45} \\
p_{15} & p_{25} & p_{35} & p_{45} & p_{55} \\
\end{pmatrix}.
\]
Then there exist unique matrices $B$ and $C$ such that $A=BCB^{T}$, where 
\[
BCB^{T}=\begin{pmatrix}
1 & 0 & 0 & 0 & 0 \\
0 & 1 & 0 & 0 & 0 \\
0 & 0 & 1 & 0 & 0 \\
b_{14} & b_{24} & b_{34} & 1 & 0 \\
b_{15} & b_{25} & b_{35} & 0 & 1 \\
\end{pmatrix}
\begin{pmatrix}
0 & 1 & 1 & 0 & 0 \\
1 & 0 & 1 & 0 & 0 \\
1 & 1 & 0 & 0 & 0\\
0 & 0 & 0 & c_{44} & c_{45} \\
0 & 0 & 0 & c_{45} & c_{55} \\
\end{pmatrix}
\begin{pmatrix}
1 & 0 & 0 & b_{14} & b_{15} \\
0 & 1 & 0 & b_{24} & b_{25} \\
0 & 0 & 1 & b_{34} & b_{35} \\
0 & 0 & 0 & 1 & 0 \\
0 & 0 & 0 & 0 & 1 \\
\end{pmatrix}.
\]
\end{lemma}
\begin{proof}
For any $p_{14}, p_{15}, p_{24}, p_{25}, p_{34}, p_{35}$, there exist unique  $b_{14}, b_{15}, b_{24}, b_{25}, b_{34}, b_{35}$ such that
%\begin{equation}\label{eq:small matrix}
\[
\begin{pmatrix}
b_{14} & b_{24} & b_{34}\\
b_{15} & b_{25} & b_{35}\\
\end{pmatrix}
\begin{pmatrix}
0 & 1 & 1 \\
1 & 0 & 1 \\
1 & 1 & 0 \\
\end{pmatrix}
=\begin{pmatrix}
-p_{14} & -p_{24} & -p_{34} \\
-p_{15} & -p_{25} & -p_{35} \\
\end{pmatrix}.
\]
This defines the invertible lower triangular matrix $B$.
The matrix $C$ is uniquely determined by the condition $C=B^{-1}A(B^{T})^{-1}$, and it is straightforward to check that $C$ has the required block diagonal shape.
%\[
%C \coloneq B^{-1}A(B^{T})^{-1}=
%\begin{pmatrix}
%0 & 1 & 1 & 0 & 0 \\
%1 & 0 & 1 & 0 & 0 \\
%1 & 1 & 0 & 0 & 0\\
%0 & 0 & 0 & c_{44} & c_{45} \\
%0 & 0 & 0 & c_{45} & c_{55} \\
%\end{pmatrix}. \qedhere
%\]
%for some $c_{44}, c_{45}, c_{55}$, which proves the existence. For the uniqueness, the given form of $BCB^T$ in the lemma implies that $b_{14}, \dots, b_{35}$ must satisfy equation \eqref{eq:small matrix}. Hence, $B$ is uniquely determined. Then $C=B^{-1}A(B^{T})^{-1}$ is also uniquely determined. 
%We define for $i\neq j$ an elementary $5\times 5$ matrix $E_{ij}(\lambda)$ with  ones on the  diagonal,  $\lambda$ at position $(i,j)$, and zeros everywhere else. Then for a $5\times 5$ matrix $X$, the product $E_{ij}(\lambda) \cdot X$ corresponds to the row operation that adds $\lambda$ times row $j$ to row $i$ of $X$. Similarly, the product $X \cdot E_{ij}(\lambda)^T = X \cdot E_{ji}(\lambda)$ corresponds to the column operation that adds $\lambda$ times column $j$ to column $i$ of $X$.
%We notice that
%\[E_{43}\left(\frac{-p_{14}}{2}\right)\cdot E_{42}\left(\frac{-p_{14}}{2}\right) \cdot E_{41}\left(\frac{p_{14}}{2}\right) \cdot A=\begin{pmatrix}
%0 & 1 & 1 & p_{14} & p_{15} \\
%1 & 0 & 1 & p_{24} & p_{25} \\
%1 & 1 & 0 & p_{34} & p_{35} \\
%0 & p_{24} & p_{34} & * & * \\
%p_{15} & p_{25} & p_{35} & p_{45} & p_{55} \\
%\end{pmatrix}\]
%One can eliminate the symmetric $p_{14}$ entry of $A$ by multiplying $A$ on the right by the  transpose of those three elementary matrices above in reversed order. In a similar way, we can eliminate the remaining entries $p_{24},~p_{34},~p_{15},~p_{25},~p_{35}$. Notice that product of several $E_{ij}$ with $i>j$ is a unit lower triangular matrix, and product of several $E_{ij}$ with $i<j$ is a unit upper triangular matrix. We recall that $\left(XY\right)^{T}=Y^TX^T$ for any $X,Y$ and thus $A=BCB^{T}$ as in the claim of the lemma.
\end{proof}

\begin{lemma}\label{lem:lor_plus_nsd}
Consider the symmetric $2 \times 2$ matrices
\[
M_1=\begin{pmatrix}
    \mu&\nu\\
    \nu&\xi
\end{pmatrix} \ \ \text{and} \ \ 
M_2=\begin{pmatrix}
\alpha&\beta\\
\beta&\gamma
\end{pmatrix}.
\]
If $M_1$ is Lorentzian, $M_2$ is negative semidefinite, and $M_1+M_2$ has nonnegative entries, then
\[
(\mu+\alpha)(\nu+\beta)(\xi+\gamma) \le \mu\nu\xi .
\]
\end{lemma}
\begin{proof}
It suffices to show that $(\mu+t\alpha)(\nu+t\beta)(\xi+t\gamma)$ is a decreasing function for $0\leq t\leq 1$. 
Since $M_1+tM_2$is a Lorentzian matrix for any $0\leq t\leq 1$, it suffices to show that  
%\begin{equation}\label{eq_1}
\[
\frac{d}{dt}(\mu+t\alpha)(\nu+t\beta)(\xi+t\gamma) \ \bigg|_{t=0}\leq 0.
\]
%\end{equation}
By taking limits, we may reduce to the case when $M_2$ is negative definite.
Since the claimed inequality is invariant under the symmetric scaling of rows and columns, we may further suppose that $\alpha=\gamma=-1$.
Then $\beta \le 1$, and hence
\[
\frac{d}{dt}(\mu+t\alpha)(\nu+t\beta)(\xi+t\gamma)\ \bigg|_{t=0}=\mu\nu\gamma+\mu\xi\beta+\nu\xi\alpha
\leq -\mu\nu+\mu\xi-\nu\xi.
\]
Since $M_1$ is Lorentzian, $\mu\xi\leq \nu^2$, so the right-hand side satisfies
\[
-\mu\nu+\mu\xi-\nu\xi \le -2\nu \sqrt{\mu\xi}+\mu\xi \le - \mu\xi \le 0.\qedhere
\]
\end{proof}

\begin{lemma}\label{lem:hard_lemma}
For positive numbers $x_1,x_2,x_3,y_1,y_2,y_3$, set
\begin{align*}
X&\coloneq 2x_1x_2+2x_1x_3+2x_2x_3-x_1^2-x_2^2-x_3^2,\\
Y&\coloneq 2y_1y_2+2y_1y_3+2y_2y_3-y_1^2-y_2^2-y_3^2,\\
Z&\coloneq x_1(y_2+y_3-y_1)+x_2(y_1+y_3-y_2)+x_3(y_1+y_2-y_3).
\end{align*}
If $X,Y,Z$ are nonnegative, then $XYZ < 32x_1x_2x_3y_1y_2y_3$.
\end{lemma}

We note that the final inequality may fail without the assumption that $X,Y,Z$ are nonnegative, for example, when $(x_1,x_2,x_3,y_1,y_2,y_3)=(6,1,1,1,1,6)$.

\begin{proof}
First suppose that $(x_1,x_2,x_3)$ and $(y_1,y_2,y_3)$ satisfy the triangle inequalities.
Setting $u_i\coloneq x_j+x_k-x_i>0$ and $v_i\coloneq y_j+y_k-y_i>0$, we have
\[
X=u_1u_2+u_1u_3+u_2u_3, \ \ Y=v_1v_2+v_1v_3+v_2v_3, \ \ 2Z=(u_2+u_3)v_1+(u_1+u_3)v_2+(u_1+u_2)v_3.
\]
The goal is to show that
\begin{multline*}
(u_1u_2+u_1u_3+u_2u_3)(v_1v_2+v_1v_3+v_2v_3)((u_2+u_3)v_1+(u_1+u_3)v_2+(u_1+u_2)v_3) \\
< (u_1+u_2)(u_1+u_3)(u_2+u_3)(v_1+v_2)(v_1+v_3)(v_2+v_3).
\end{multline*}
The right-hand side minus the left-hand side can be simplified to
\begin{multline*}
u_1^2 u_2 v_1^2 v_2+ u_1^2 u_3 v_1^2 v_2+ u_1^2 u_2 v_1^2 v_3+ u_1^2 u_3 v_1^2 v_3+ u_1 u_2^2 v_1 v_2^2  + u_2^2 u_3 v_1 v_2^2  + u_1 u_2^2 v_2^2 v_3 + u_2^2 u_3 v_2^2 v_3\\
 + u_1 u_3^2 v_2 v_3^2 + u_2 u_3^2 v_2 v_3^2  + u_1 u_3^2 v_1 v_3^2  + u_2 u_3^2 v_1 v_3^2 - 2 u_1 u_2 u_3 v_1 v_2 v_3.
\end{multline*}
Notice that at most one of the nonnegative numbers $u_1,u_2,u_3$ is zero, and similarly for $v_1,v_2,v_3$. Thus, the sum of the positive coefficient terms is strictly positive. If any of the $u_i$ is zero, then the difference is positive. If $u_1,u_2,u_3$ are all positive, then
the only negative term $- 2 u_1 u_2 u_3 v_1 v_2 v_3 $ can be grouped with the terms $u_1^2 u_2 v_1^2 v_2$ and $u_2 u_3^2 v_2 v_3^2$ to form a square, and the remaining sum is positive.

Thus, without loss of generality, it is sufficient to justify the inequality under the assumption that  $x_1>x_2+x_3$ and $x_2 \ge x_3$. If $y_3<y_2$, then $(y_1+y_3-y_2)<(y_1+y_2-y_3)$. In this case, 
\[
x_2(y_1+y_3-y_2)+x_3(y_1+y_2-y_3)
<
x_2(y_1+y_2-y_3)+x_3(y_1+y_3-y_2)  \ \ \text{and} \ \ 
\]
\[
x_1(y_2+y_3-y_1)+x_2(y_1+y_3-y_2)+x_3(y_1+y_2-y_3)
<
x_1(y_2+y_3-y_1)+x_2(y_1+y_2-y_3)+x_3(y_1+y_3-y_2).
\]
%This follows from the sequence of equivalent inequalities
%\begin{align*}
%x_2(y_1+y_3-y_2)+x_3(y_1+y_2-y_3)
%&<
%x_2(y_1+y_2-y_3)+x_3(y_1+y_3-y_2),\\
%(x_2-x_3)(y_1+y_3-y_2)
%&<
%(x_2-x_3)(y_1+y_2-y_3),\\
%(y_1+y_3-y_2)
%&<
%(y_1+y_2-y_3),\\
%2y_3
%&<
%2y_2.
%\end{align*}
Thus,  swapping $y_2$ and $y_3$ makes $Z$ larger without changing $X$, $Y$, and $x_1x_2x_3y_1y_2y_3$.
Therefore, it is enough to consider the case when $x_1>x_2+x_3$ and $x_2 \ge x_3$ and $y_3 \ge y_2$.

Since
$X= -(x_1-x_2-x_3)^2+4x_2x_3$,
 replacing $x_1$ by $x_2+x_3$ makes $X$ larger without changing $Y$. We check that the same replacement makes the ratio $Z/x_1$ larger:
\begin{multline*}
(y_2+y_3-y_1)+(x_2(y_1+y_3-y_2)+x_3(y_1+y_2-y_3))x_1^{-1}\\
< (y_2+y_3-y_1)+(x_2(y_1+y_3-y_2)+x_3(y_1+y_2-y_3))(x_2+x_3)^{-1}.
\end{multline*}
%In fact, since $x_1>x_2+x_3$, the above inequality is equivalent to 
%\[
% x_3(y_3-y_1-y_2) < x_2(y_1+y_3-y_2),
%\]
%which holds because $x_2 \ge x_3>0$ and $y_1+y_3-y_2>y_3-y_1-y_2$.
In fact, we have $x_1>x_2+x_3$ and 
\[
x_2(y_1+y_3-y_2)+x_3(y_1+y_2-y_3)=(x_2+x_3)y_1+(x_2-x_3)(y_3-y_2)\geq (x_2+x_3)y_1> 0.
\]
Since
\[
Z=y_1(x_2+x_3-x_1)+y_2(x_1+x_3-x_2)+y_3(x_1+x_2-x_3),
\]
replacing $x_1$ by $x_2+x_3$ preserves the nonnegativity of $Z$, this reduces the problem to the case when $(x_1,x_2,x_3)$ satisfies the triangle inequalities. Repeating the argument for $(y_1, y_2, y_3)$, we reduce the problem to the case where both $(x_1, x_2, x_3)$ and $(y_1, y_2, y_3)$ satisfy the triangle inequalities, which was considered before.
\end{proof}

%We are ready to prove Theorem \ref{mth:pentagonal}.
\begin{proof}[Proof of Theorem~\ref{mth:pentagonal}]
%By Theorem \ref{thm:BR}, $\BR'(\underline{\L_5^+})$ is dual to $\Cut_5$. Hence, we know that the set of primitive ratios are given by the following ratios, up to permutations of indices \cite[Section 30.6]{Deza-Laurent}:
%	\begin{enumerate}
%		\item $\dfrac{p_{11}p_{23}}{p_{12}p_{13}}$,
 %       \smallskip
	%	\item $\dfrac{p_{12}p_{13}p_{23}p_{44}p_{45}p_{55}}{p_{14}p_{24}p_{34}p_{15}p_{25}p_{35}}$. 
	%\end{enumerate}
%Let $(p_{ij})_{ij}\in \L_5^+$. In \eqref{eq:triangular}, we have proved first ratio is bounded by 2. To see that 2 is the optimal constant, evaluate this ratio on the family of Lorentzian matrices  $\left(\begin{smallmatrix}
%	2 & 1+x & 1+x \\
	%1+x & x & 1 \\
	%1+x & 1 & x
%\end{smallmatrix}\right)$ where $0<x<1$ as $x$ approaches 0.
%Even though it follows from Theorem \ref{thm:BR} that the pentagonal ratio is bounded, we will prove the second part by showing the pentagonal ratio is at most 4, and showing that 4 is the optimal ratio. 
%The remaining part of the theorem is to show that the pentagonal ratio is bounded, and the optimal  constant is $4$.  We start by showing that the optimal  constant is at most 4.    

To show the pentagonal ratio is at most 4, it suffices to consider Lorentzian matrices of the form %\textcolor{blue}{(Why? Is it Corollary 3.5. in ``Lorentzian polynomials'': The multi-affine part of any Lorentzian polynomial is a Lorentzian polynomial?)}
\[
A=\begin{pmatrix}
0 & 1 & 1 & p_{14} & p_{15} \\
1 & 0 & 1 & p_{24} & p_{25} \\
1 & 1 & 0 & p_{34} & p_{35} \\
p_{14} & p_{24} & p_{34} & p_{44} & p_{45} \\
p_{15} & p_{25} & p_{35} & p_{45} & p_{55} \\
\end{pmatrix}.
\]
Indeed, we can set the first three diagonal entries to be zero, since decreasing diagonal entries only decreases eigenvalues. %and hence preserving Lorentzian matrices. %Clearly, the pentagonal ratio is defined on such matrices. 
We can then rescale rows and columns using the scaling action to obtain matrix $A$. Both operations do not affect the value of the pentagonal ratio.
We need to show that 
\[
\frac{p_{44}p_{45}p_{55}}{p_{14}p_{24}p_{34}p_{15}p_{25}p_{35}}\leq 4.
\]
By Lemma \ref{lem:factors}, there is a unique factorization 
\[
A=B \cdot C \cdot B^{T}=\begin{pmatrix}
1 & 0 & 0 & 0 & 0 \\
0 & 1 & 0 & 0 & 0 \\
0 & 0 & 1 & 0 & 0 \\
b_{14} & b_{24} & b_{34} & 1 & 0 \\
b_{15} & b_{25} & b_{35} & 0 & 1 \\
\end{pmatrix}
\begin{pmatrix}
0 & 1 & 1 & 0 & 0 \\
1 & 0 & 1 & 0 & 0 \\
1 & 1 & 0 & 0 & 0\\
0 & 0 & 0 & c_{44} & c_{45} \\
0 & 0 & 0 & c_{45} & c_{55} \\
\end{pmatrix}
\begin{pmatrix}
1 & 0 & 0 & b_{14} & b_{15} \\
0 & 1 & 0 & b_{24} & b_{25} \\
0 & 0 & 1 & b_{34} & b_{35} \\
0 & 0 & 0 & 1 & 0 \\
0 & 0 & 0 & 0 & 1 \\
\end{pmatrix}.
\]
Since $A$ is Lorentzian, $C$ has at most one positive eigenvalue. The upper-left block of $C$ has exactly one positive eigenvalue, it follows that the lower-right block of $C$ is negative semidefinite.
%left upper block of $C$
%\[
%bl_1=\begin{pmatrix}
%0 & 1 & 1 \\
%1 & 0 & 1\\
%1 & 1 & 0\\    
%\end{pmatrix}
%\]
%has one positive eigenvalue, and thus the lower right block of $C$
%\[
%bl_2=\begin{pmatrix}
%c_{44} & c_{45}\\
%c_{45} & c_{55}\\    
%\end{pmatrix}
%\]
%is negative semidefinite. 
For $0 \le t \le 1$, we consider the matrices $A(t)$ and $C(t)$ defined by
%Let $C(t)$ be a matrix obtained from $C$ by rescaling $bl_2$ by parameter $0 \leq t \leq 1$. We define $A(t)$ as follows
\[
A(t)=B\cdot C(t)\cdot B^{T}=\begin{pmatrix}
1 & 0 & 0 & 0 & 0 \\
0 & 1 & 0 & 0 & 0 \\
0 & 0 & 1 & 0 & 0 \\
b_{14} & b_{24} & b_{34} & 1 & 0 \\
b_{15} & b_{25} & b_{35} & 0 & 1 \\
\end{pmatrix}
\begin{pmatrix}
0 & 1 & 1 & 0 & 0 \\
1 & 0 & 1 & 0 & 0 \\
1 & 1 & 0 & 0 & 0\\
0 & 0 & 0 & tc_{44} & tc_{45} \\
0 & 0 & 0 & tc_{45} & tc_{55} \\
\end{pmatrix}
\begin{pmatrix}
1 & 0 & 0 & b_{14} & b_{15} \\
0 & 1 & 0 & b_{24} & b_{25} \\
0 & 0 & 1 & b_{34} & b_{35} \\
0 & 0 & 0 & 1 & 0 \\
0 & 0 & 0 & 0 & 1 \\
\end{pmatrix}.
\]
Matrices $B$ and $B^{T}$ correspond to sequences of row and column operations respectively, so 
\[
A(t)=
%\begin{pmatrix}
%0 & 1 & 1 & p_{14} & p_{15} \\
%1 & 0 & 1 & p_{24} & p_{25} \\
%1 & 1 & 0 & p_{34} & p_{35} \\
%p_{14} & p_{24} & p_{34} & p_{44}-(1-t)c_{44} & p_{45}-(1-t)c_{45} \\
%p_{15} & p_{25} & p_{35} & p_{45}-(1-t)c_{45} & p_{55}-(1-t)c_{55} \\
%\end{pmatrix}=
\begin{pmatrix}
0 & 1 & 1 & p_{14} & p_{15} \\
1 & 0 & 1 & p_{24} & p_{25} \\
1 & 1 & 0 & p_{34} & p_{35} \\
p_{14} & p_{24} & p_{34} & X+tc_{44} & Y+tc_{45} \\
p_{15} & p_{25} & p_{35} & Y+tc_{45} & Z+tc_{55} \\
\end{pmatrix}
%; \;\text{in particular} \;
%A(0)=\begin{pmatrix}
%0 & 1 & 1 & p_{14} & p_{15} \\
%1 & 0 & 1 & p_{24} & p_{25} \\
%1 & 1 & 0 & p_{34} & p_{35} \\
%p_{14} & p_{24} & p_{34} & X & Y \\
%p_{15} & p_{25} & p_{35} & Y & Z \\
%\end{pmatrix}.
\]
for some $X,Y,Z$ independent of $t$. 
%We have
%\begin{equation}
%\label{eq:A0}
% A(0)=\begin{pmatrix}
%0 & 1 & 1 & p_{14} & p_{15} \\
%1 & 0 & 1 & p_{24} & p_{25} \\
%1 & 1 & 0 & p_{34} & p_{35} \\
%p_{14} & p_{24} & p_{34} & X & Y \\
%p_{15} & p_{25} & p_{35} & Y & Z \\
%\end{pmatrix}
%\ \ \text{and}\ \  A(1)=\begin{pmatrix}
%0 & 1 & 1 & p_{14} & p_{15} \\
%1 & 0 & 1 & p_{24} & p_{25} \\
%1 & 1 & 0 & p_{34} & p_{35} \\
%p_{14} & p_{24} & p_{34} & p_{44} & p_{45} \\
%p_{15} & p_{25} & p_{35} & p_{45} & p_{55} \\
%\end{pmatrix}.
%\end{equation}
%Furthermore, $p_{44}=X+c_{44}$, $p_{45}=Y+c_{45}$ and $p_{55}=Z+c_{55}$.
%Since the rank of a product of matrices is the minimum of the rank of each factor, we have  
The rank of $C(0)$ is $3$, and hence the rank of $A(0)$ is $3$. Thus, all the $4\times 4$ minors of $A(0)$ are zero, which implies that 
\begin{align*}
    X&=\frac{2p_{14}p_{24}+2p_{14}p_{34}+2p_{24}p_{34}-p_{14}^2-p_{24}^2-p_{34}^2}{2},\\
    Z&=\frac{2p_{15}p_{25}+2p_{15}p_{35}+2p_{25}p_{35}-p_{15}^2-p_{25}^2-p_{35}^2}{2}, \\
Y&=\frac{p_{14}p_{25}+p_{14}p_{35}+p_{24}p_{15}+p_{24}p_{35}+p_{34}p_{15}+p_{34}p_{25}-p_{14}p_{15}-p_{24}p_{25}-p_{34}p_{35}}{2}.
\end{align*}
Since $c_{44}$ and $c_{55}$ are nonpositive and $p_{44}=X+c_{44}$ and  $p_{55}=Z+c_{55}$ are nonnegative, 
\[
X+tc_{44} \ge 0 \ \ \text{and} \ \ Z+tc_{55} \ge 0 \ \ \text{for all $0 \le t \le 1$.}
\]
In particular, $X$ and $Z$ are nonnegative.

We now reduce to the case when $Y$ is nonnegative as well.
Since $C(t)$ has at most one positive eigenvalue for all $0 \le t \le 1$,  the same holds for the matrix $A(t)$, and hence
\[
\det \begin{pmatrix}
X+tc_{44} & Y+tc_{45} \\
 Y+tc_{45} & Z+tc_{55} \\    
\end{pmatrix} \le 0.
\]
Suppose $Y$ is negative. Since $Y+c_{45} \ge 0$, we must have $c_{45}>0$ and there is a unique positive $t_1 \le 1$ such that $Y+t_1c_{45}=0$. The determinantal inequality above says that 
\[
X+t_1c_{44}=0  \ \ \text{or} \ \  Z+t_1c_{55}=0.
\]
If $t_1<1$, this implies that $c_{44}=0$ or $c_{55}=0$, and hence $c_{45}=0$ by the negative semidefiniteness of the lower-right block of $C$, reaching a contradiction.
Thus 
\[
p_{45}=Y+c_{45}=0,
\]
in which case the pentagonal inequality holds trivially.

The remaining case to consider is when $X,Y,Z \ge 0$, that is, when the matrix $A(0)$ is Lorentzian. 
We apply Lemma \ref{lem:lor_plus_nsd} to the sum
\[
\begin{pmatrix}
X & Y \\
 Y & Z \\    
\end{pmatrix}+\begin{pmatrix}
c_{44} & c_{45} \\
c_{45} & c_{55} \\    
\end{pmatrix}=\begin{pmatrix}
p_{44} & p_{45} \\
p_{45} & p_{55} \\    
\end{pmatrix},
\]
and conclude that $p_{44}p_{45}p_{55}\leq XYZ$. Thus, the pentagonal ratio of $A(0)$ is greater or equal to the pentagonal ratio of $A$. Thus, it is enough to show the pentagonal inequality for $A(0)$, which reads
\[
XYZ\leq 4p_{14}p_{24}p_{34}p_{15}p_{25}p_{35}.
\]
This is precisely the content of Lemma \ref{lem:hard_lemma}.
%It remains to show that the constant 4 is optimal.
%Let $x$ be a positive number, and consider the matrix
%\[
%M=
%\begin{pmatrix}
%0 & 1 & 1 & x & 2 + x \\[1pt]
%1 & 0 & 1 & 2 & 2 \\[1pt]
%1 & 1 & 0 & 2 + x & x \\[1pt]
%x & 2 & 2 + x & 4x & 4 + 4x \\[1pt]
%2 + x & 2 & x & 4 + 4x & 4x
%\end{pmatrix}.
%\]
%Notice that $M$ is of rank 3. Since the first $3\times 3$ minor of $M$ is a nondegenerate Lorentzian matrix, $M$ is also a Lorentzian matrix for all $x\ge 0$. 
%Then $M$ is in the form of $A(0)$ as in \eqref{eq:A0}, satisfying \eqref{eq:xyz}, so it is a Lorentzian matrix for all  $x \ge 0$. 
%We have
%\[
%\lim_{x \to 0} \frac{p_{12}p_{13}p_{23}p_{45}p_{44}p_{55}}{p_{14}p_{15}p_{24}p_{25}p_{34}p_{35}}=\lim_{x \to 0} \frac{(4x)^2(4+4x)}{4x^2(2+x)^2}=4. 
%\]
%Since the pentagonal ratio continuously extends to Lorentzian matrices with $p_{11}=p_{22}=p_{33}=0$ and the rest of the entries strictly positive, it follows that there is a family in $\L_5^+$ on which the evaluation of the pentagonal ratio limits to 4.
\end{proof}

It is interesting to compare the above analysis with that for the pentagonal ratio on $\Delta_n^+(\T_p)$, which is much simpler.

\begin{proposition}
Let nonnegative $p$ and any matrix $(p_{ij}) \in \Delta_n^+(\T_p)$, we have
\[
\frac{p_{12}p_{13}p_{23}p_{45}p_{44}p_{55}}{p_{14}p_{15}p_{24}p_{25}p_{34}p_{35}} \le 8^{p}.
\]
The equality is achieved by the matrix 
\[
\begin{pmatrix}
0 & 2^{\frac{p}{2}} & 2^{\frac{p}{2}} & 1 & 1 \\
2^{\frac{p}{2}} & 0 & 2^{\frac{p}{2}} & 1 & 1 \\
2^{\frac{p}{2}} & 2^{\frac{p}{2}} & 0 & 1 & 1 \\
1 & 1 & 1 & 2^{\frac{p}{2}} & 2^{\frac{p}{2}} \\
1 & 1 & 1 & 2^{\frac{p}{2}} & 2^{\frac{p}{2}}
\end{pmatrix} \in \Delta_n(\T_p).
\]
\end{proposition}
\begin{proof}
We first consider the case when $p$ is positive. We have
\[
(p_{44}p_{13})^{1/p} \le 2 (p_{14}p_{34})^{1/p}, ~~ (p_{55}p_{23})^{1/p} \le 2 (p_{25}p_{35})^{1/p}, \ \ \text{so} \ \ (p_{44}p_{55}p_{13}p_{23})^{1/p} \le 4 (p_{14}p_{34}p_{25}p_{35})^{1/p}.
\]
Similarly, we have
\[
(p_{44}p_{23})^{1/p} \le 2 (p_{24}p_{34})^{1/p}, ~~ (p_{55}p_{13})^{1/p} \le 2 (p_{15}p_{35})^{1/p}, \ \ \text{so} \ \ (p_{44}p_{55}p_{13}p_{23})^{1/p} \le 4 (p_{24}p_{34}p_{15}p_{35})^{1/p}.
\]
Using the triangle inequality $(p_{12}p_{45})^{1/p} \le (p_{14}p_{25})^{1/p}+(p_{15}p_{24})^{1/p}$, we get
\begin{align*}
(p_{12}p_{13}p_{23}p_{44}p_{55}p_{45})^{1/p} &\le (p_{44}p_{55}p_{13}p_{23}p_{14}p_{25})^{1/p}+(p_{44}p_{55}p_{13}p_{23}p_{15}p_{24})^{1/p} \\
&\le 8(p_{14}p_{15}p_{24}p_{25}p_{34}p_{35})^{1/p}.
\end{align*}
%Raising both sides to the $p$-th power yields 
%\[
%\frac{p_{12}p_{13}p_{23}p_{45}p_{44}p_{55}}{p_{14}p_{15}p_{24}p_{25}p_{34}p_{35}} \le 2^{3p}.
%\]
Since $\Delta_n(\T_0) \subseteq \Delta_n(\T_p)$ for all $p>0$, the case of $p=0$ follows from taking the limit $p \to 0$.
\end{proof}


\section{Lorentzian matrices and $\delta$-hyperbolic spaces}\label{sec:hyperbolic}

%Related background.

%Proof of Theorem B.

We give a detailed proof of Theorem~\ref{mth:BR} outlined in the introduction.


\begin{proposition}
\label{prop:inclusions}
	For any $n>0$, we have 
    $\Delta_n(\T_0)\subseteq \L_n\subseteq \Delta_n(\T_2)$.
\end{proposition}

In the language of \cite{BHKL-polymatroids}, the proposition states that every rank $2$ polymatroid over $\T_0$ is Lorentzian, and every Lorentzian matrix is a rank $2$ polymatroid over $\T_2$.

\begin{proof}
The first inclusion is a special case of  \cite[Corollary 3.16]{Branden-Huh} for quadratic polynomials. The second inclusion follows from \cite{BHKL-triangular2}, whose proof we reproduce here.

Suppose $(p_{ij})$ is a Lorentzian matrix. If $i,j,k,l$ are distinct indices, we consider the matrix 
\[
\begin{pmatrix}
0 & p_{ij} & p_{ik} & p_{il} \\
p_{ij} & 0 & p_{jk} & p_{jl}  \\
p_{ik} & p_{jk} & 0 & p_{kl}  \\
p_{il} & p_{jl} & p_{kl} & 0
\end{pmatrix} 
\]
obtained from replacing the diagonal entries of a principal submatrix of $(p_{ij})$ by zero. This matrix is Lorentzian with 
\begin{multline*}
\det=-\left(\sqrt{p_{ij}p_{kl}}+\sqrt{p_{ik}p_{jl}}+\sqrt{p_{il}p_{jk}}\right)\cdot(-\sqrt{p_{ij}p_{kl}}+\sqrt{p_{ik}p_{jl}}+\sqrt{p_{il}p_{jk}})\\
      \cdot(\sqrt{p_{ij}p_{jl}}-\sqrt{p_{ik}p_{jl}}+\sqrt{p_{il}p_{jk}})\cdot(\sqrt{p_{ij}p_{kl}}+\sqrt{p_{ik}p_{jl}}-\sqrt{p_{il}p_{jk}}) \le 0,
\end{multline*}
so $\sqrt{p_{ij}p_{kl}}$, $\sqrt{p_{ik}p_{jl}}$, $\sqrt{p_{il}p_{jk}}$ form three sides of a triangle \cite[Proposition 3.1]{HHMWW}.

If there are exactly three distinct indices among $i,j,k,l$, we may suppose $k=l$. 
We need to show that $\sqrt{p_{ij}p_{kk}}$, $\sqrt{p_{ik}p_{jk}}$,
 $\sqrt{p_{ik}p_{jk}}$ form sides of a triangle, that is, $\sqrt{p_{ij}p_{kk}} \le 2\sqrt{p_{ik}p_{jk}}$.
This follows from the Lorentzian condition for $(p_{ij})$ because
\[
\det\begin{pmatrix}
	0 & p_{ij} & p_{ik} \\
	p_{ij} & 0      & p_{jk} \\
	p_{ik} & p_{jk} & p_{kk}
\end{pmatrix}=p_{ij}(2p_{ik}p_{jk}-p_{ij}p_{kk})\ge 0.
\]
If three of $i,j,k,l$ are equal, then the statement is trivial. Thus, the only remaining nontrivial case is when $i=k$ and $j=l$. We need to check that $\sqrt{p_{ii}p_{jj}} \le 2p_{ij}$, and this follows from
\[
\det\begin{pmatrix}
	p_{ii} & p_{ij}  \\
	p_{ij} & p_{jj}      
\end{pmatrix}=p_{ii}p_{jj}-p_{ij}^2\le 0. \qedhere
\]
\end{proof}

We now recall the definition of tree metric in terms of phylogenetic trees, following \cite[Section 4.3]{Maclagan-Sturmfels}. A \emph{phylogenetic} tree on $n$ leaves is a tree with $n$ labeled leaves and no vertices of degree 2.
We say that $(d_{ij})_{1 \le i< j \le n}$ is a \emph{tree metric} if there is a phylogenetic tree $\tau$ with $n$ leaves and nonnegative edge lengths $\ell$ on $\tau$ such that
\[
\mathrm{d}_{ij}=\bigl( \text{the sum of all $\ell(e)$ over all edges $e$ in the unique path from $i$ to $j$ in $\tau$}\bigr).\footnote{In \cite[Section 4.3]{Maclagan-Sturmfels}, edge lengths are allowed to be non-positive. A vector $d$ is a tree metric in this sense if and only if the sum of $d$ and a constant vector of large positive value is a tree metric in the more restrictive sense above.}
\]
Every tree metric uniquely determines a phylogenetic tree together with the assignment of edge lengths. 

\begin{lemma}
\label{lem:convhullT0} 
The following statements hold for $\underline{\Delta}_n^+(\T_0)$.
\leavevmode
\begin{enumerate}[(1)]\itemsep 5pt
    \item The set of coordinatewise logarithms $\log \underline{\Delta}_n^+(\T_0)\subseteq \RR^{\binom{n}{2}}$ is equal to the space of tree metrics on $n$ leaves. %where we consider each tree metric $d:\binom{[n]}{2}\to \RR$ as an element in $\RR^{\binom{[n]}{2}}$.
    \item 	The convex hull of $\log \underline{\Delta}_n^+(\T_0)$ in $\RR^{\binom{n}{2}}$
    is equal to the cut cone $\Cut_n$.
    \item The cone of reduced bounded ratios $\underline{\BR}(\Delta_n^+(\T_0))$ is the dual of the cut cone $\Cut_n$.
\end{enumerate}
\end{lemma}



\begin{proof}%[Proof of Lemma~\ref{lem:convhullT0}]
Note that the points in  $\log \underline{\Delta}_n^+(\T_0)$ are of the form $(q_{ij})=(\log p_{ij})$, where 
\[
q_{ii}=\log p_{ii}=0 \ \ \text{and} \ \ q_{ij} =\log p_{ij} \ge \log \sqrt{p_{ii}p_{jj}} =0,
\]
and in addition, for distinct indices $i,j,k,l$, we have
\[
q_{ij}+q_{kl}=\log p_{ij}p_{kl} \le \log \max(p_{ik}p_{jl},p_{il}p_{jk}) = \max(q_{ik}+q_{jl}, q_{il}+q_{jk}).
\]
Thus,  the first statement is the characterization of tree metrics by the ``four-point condition'', which is exactly the logarithm of the condition defining $\underline{\Delta}^+_n(\T_0)$ \cite[Lemma 4.3.6]{Maclagan-Sturmfels}.


	%By definition, $\underline{\Delta(\T_0)_n^+}$ consists of positive matrices $(p_{ij})_{ij}$ with ones on the diagonal such that for all $i,j,k,l\in [n]$, $p_{ij}p_{kl}\le \max \{p_{ik}p_{jl}, p_{il}p_{jk}\}$. Let $q_{ij}=\log p_{ij}$. Then
%\begin{multline}
%\label{eq:logT0}
%	\log \underline{\Delta(\T_0)_n^+}= \Bigl\{
%	(q_{ij})_{ij}\in \Sym_n(\RR)\;\Bigm|\; q_{ii}=0 \text{ for all } i\in [n],\\
%	\text{ and for all }i,j,k,l\in [n],\; 
%	q_{ij}+q_{kl}\le \max \{q_{ik}+q_{jl}, q_{il}+q_{jk}\}
%	\Bigr\}
%\end{multline}
%By the first inclusion in Proposition~\ref{prop:inclusions}, $(p_{ij})_{ij}$ is Lorentzian, and therefore all the off diagonal entries $p_{ij}$ for $i\neq j$ must be at least 1. This means that in fact $\log \underline{\Delta(\T_0)_n^+}$ contains only matrices with nonnegative entries. 
%The conditions in \eqref{eq:logT0} are exactly the ``four-point condition'' on tree metrics, and by \cite[Lemma 4.3.6]{Maclagan-Sturmfels}, we see that $\log \underline{\Delta(\T_0)_n^+}$ is  the space of  tree metrics on $n$ vertices. Thus, we have proved statement (i).

For the second statement, observe that the cut vector
 $\delta(S)$ is the tree metric corresponding to the phylogenetic tree that consists of a single edge of length $1$ that connects the leaves in $S$ with the leaves not in $S$. 
Therefore, the extremal rays of $\Cut_n$ are in $\log \underline{\Delta}_n^+(\T_0)$. 
For the other inclusion, consider a tree metric $d$  given by a phylogenetic tree $\tau$ and a length function $\ell$. For each edge $e$ of $\tau$, let $S(e)$ be the set of leaves that remain connected to the distinguished leaf $1$ after removing $e$ from $\tau$. Then $d$ admits the decomposition
\[
d=\sum_e \ell(e) \delta(S(e)),
\]
showing that $d$ is in the cut cone.
This proves the second statement.


For the third statement, notice that
\begin{align*}
\underline{\BR}\left(\Delta_n^+(\T_0)\right)
&=\left\{ \alpha \in \RR^{\binom{n}{2}}  \;\middle|\;  \exists c> 0\text{ such that } \sum_{1\le i< j\le n}\alpha_{ij}q_{ij}\le 	\log c \ \ \text{for all $q\in \log \underline{\Delta}_n^+(\T_0)$}\right\}\\
&=\left\{ \alpha\in \RR^{\binom{n}{2}}  \;\middle|\;  \exists c> 0\text{ such that } \sum_{1\le i< j\le n}\alpha_{ij}q_{ij}\le 	\log c \ \ \text{for all $q\in \Cut_n$}\right\}\\
&= \left\{ \alpha \in \RR^{\binom{n}{2}}  \;\middle|\;  \sum_{1\le i< j\le n}\alpha_{ij}q_{ij}\le 0 \ \ \text{for all $q\in \Cut_n$}\right\}.
\end{align*}
The first equality is the definition of $\underline{\BR}$, the second follows from the fact that $\Cut_n$ is the convex hull of $\log \underline{\Delta}_n^+(\T_0)$, and the third follows from the fact that $\Cut_n$ is a cone.
\end{proof}



We now review Gromov's tree approximation theorem for $\delta$-hyperbolic spaces. %\footnote{The idea of using Gromov’s theorem on $\delta$-hyperbolic spaces in connection with Lorentzian polynomials originated with the manuscript under preparation [BHKL2].}  
Let $\delta$ be a nonnegative real number. A metric $d$  on $[n]$ 
is called \emph{$\delta$-hyperbolic} if
\[
d_{ij}+d_{kl}\le \max \bigl(d_{ik}+d_{jl}+2\delta, d_{il}+d_{jk}+2\delta\bigr) \ \ \text{for all $i,j,k,l\in [n]$.}
\]
As noted above, when $\delta=0$, the condition is equivalent to $d$ being a tree metric. 
We will use the following corollary of the tree approximation theorem; see {\cite[Section 6.1]{Gromov} and \cite[Chapitre 2]{GdlH}}.

\begin{proposition}%[{\cite[Section 6.1]{Gromov} and \cite[Chapitre 2]{GdlH}}]
	\label{prop:tree-approx}
	For any positive real number $\delta$ and a positive integer $n$,  
there is a constant $c$ %depending only on $\delta$ and $n$, 
such that any $\delta$-hyperbolic metric $d$ on $[n]$ can be approximated by a $0$-hyperbolic metric $d'$: 
		\[
        d_{ij}-c\le d'_{ij}\le d_{ij} \ \ \text{for all $i,j \in [n]$}.
        \]
\end{proposition}

\begin{corollary}
\label{lem:finite-hausdorff}
For any positive integer $n$, there is a constant $e^c$ such that
	for any $(p_{ij})\in \underline{\Delta}_n^+(\T_2)$, there is $(p_{ij}')\in \underline{\Delta}_n^+(\T_0)$ satisfying the condition
	\[e^{-c}p_{ij}\le p_{ij}' \le p_{ij}\ \ \text{for all }i,j\in [n].\]
\end{corollary}
\begin{proof}
	We reformulate the statement in terms of log coordinates, where $q_{ij}=\log p_{ij}$. Since $(p_{ij})\in \underline{\Delta}_n^+(\T_2)$, we have 
	\begin{enumerate}[(1)]\itemsep 5pt
	\item $q_{ii}=0$ for all $i\in [n]$;
\item $q_{ij}\leq q_{ik}+q_{jk}+2\log 2$ for all distinct $i,j,k\in [n]$;
\item $q_{ij}+q_{kl}\leq \max(q_{ik}+q_{jl}+2\log2, q_{il}+q_{jk}+2\log 2)$ for all distinct $i,j,k,l\in [n]$,
	\end{enumerate}
	where the last condition follows from 
    \[
    \sqrt{p_{ij}p_{kl}}\leq \sqrt{p_{ik}p_{jl}}+\sqrt{p_{il}p_{jk}}\leq \max(2\sqrt{p_{ik}p_{jl}}, 2\sqrt{p_{il}p_{jk}}).
    \]
Thus any such $(q_{ij})$ is a $\delta$-hyperbolic metric on $[n]$ for $\delta=\log 2$. The conclusion follows from
	Lemma~\ref{lem:convhullT0}  and Proposition~\ref{prop:tree-approx}.
    %there is a constant $c$ and some $(q'_{ij})\in \log \underline{\Delta}_n^+(\T_0)$ such that 
	%\[q_{ij}-c\le q'_{ij}\le q_{ij},\]
    %which is equivalent to the desired inequality for $C=\exp c$.
\end{proof}

For the following lemma, let $X\subseteq Y$ be arbitrary subsets of $n \times n$ symmetric matrices with positive entries.

\begin{lemma}
\label{lem:same-BR}
Suppose that there is a constant $e^c$ such that, 
for any $(p_{ij}) \in Y$, there is 
$(p'_{ij})\in X$ satisfying $e^{-c}p_{ij}\le p'_{ij}\le p_{ij}$ for all $i,j \in [n]$. Then $\BR(X)=\BR(Y)$.
\end{lemma}
\begin{proof}
Since $X\subseteq Y$, we have $\BR(Y)\subseteq \BR(X)$. 
We prove the reverse inclusion. Given $(p_{ij}) \in Y$,
our assumption says that there is $(p_{ij}') \in X$ such that
\[
e^{-c}p_{ij}\le p'_{ij}\le p_{ij} \ \ \text{for all $i,j \in [n]$.}
\]
%This implies that
%\[
%\prod_{1\leq i\le j\leq n} p_{ij}^{\alpha_{ij}}\leq \prod_{1\leq i\le j\leq n}C^{|\alpha_{ij}|} \cdot\prod_{1\leq i\le j\leq n}  {p_{ij}'}^{\alpha_{ij}}.
%\]
If $\prod_{i\le  j}  {p_{ij}'}^{\alpha_{ij}}$ is bounded above by a constant, then the same is true for $\prod_{i\le  j}  {p_{ij}}^{\alpha_{ij}}$.
%    Suppose to the contrary that $\BR(X')\supsetneq \BR(X)$ and $\alpha \in \BR(X')\setminus \BR(X)$. Then, by the definition of bounded ratios, there exists a sequence $\{p_m\}_{m>0}$ in $X$ such that 
%	\[\prod_{i\le j}\big((p_m)_{ij}\big)^{\alpha_{ij}}>m .\]
%	By assumption, for each $m$, there exists $p_m'$ such that
%	\[C^{-1}(p_m)_{ij}\le (p_m')_{ij}\le (p_m)_{ij}.\]
%	Note that 
%\begin{align*}
%    \prod_{i\le j}((p'_m)_{ij})^{\alpha_{ij}} &=\left(\prod_{\substack{i\le j\\\alpha_{ij}\geq 0}}((p'_m)_{ij})^{\alpha_{ij}}\right)\left(\prod_{\substack{i\le j\\\alpha_{ij}<0}}((p'_m)_{ij})^{\alpha_{ij}}\right) \\ 
%    &\geq
%\left(C^{-\sum_{\alpha_{ij}\geq 0}\alpha_{ij}}  \prod_{\substack{i\le j\\\alpha_{ij}\geq 0}}((p_m)_{ij})^{\alpha_{ij}}\right) \left(\prod_{\substack{i\le j\\\alpha_{ij}<0}}((p_m)_{ij})^{\alpha_{ij}}\right) \\
%&> m C^{-\sum_{\alpha_{ij}\geq 0}\alpha_{ij}}  ,
%\end{align*}
%	contradicting the assumption that
%there exists $K>0$ such that $\prod_{i\le j}(p'_{ij})^{\alpha_{ij}}<K$ for any $(p'_{ij})_{ij}\in X'$.
\end{proof}


\begin{proof}[Proof of Theorem~\ref{mth:BR}]
Proposition~\ref{prop:inclusions}, Corollary~\ref{lem:finite-hausdorff}, and 
Lemma~\ref{lem:same-BR}, we have
\[
\underline{\BR}(\Delta_n^+(\T_0))=\underline{\BR}(\L_n^+)=\underline{\BR}(\Delta_n^+(\T_2)).
\]
Lemma~\ref{lem:convhullT0} says that this cone is dual to the cut cone $\Cut_n$.
\end{proof}






\section{Optimal bounding constants for $\L_3^+$}

We determine the optimal bounding constants on $\underline{\BR}(\L_3^+)$.
%We prove Theorem~\ref{mth:entropy}.
%The characterization of multiplicative inequalities seems to be a rather complicated problem. 
%
%For $\alpha \in \BR(\L_n^+)$, we define an optimal upper bound
%\[f_{\alpha}\coloneqq \inf \left\{C_\alpha>0 \text{ such that }\forall p\in \L_n^+, \prod_{1\le i\leq j\le n}p_{ij}^{\alpha_{ij}}\le C_\alpha\right\}.\]
%Note that the correctness of the above definition follows from the fact that $\BR(\L_n^+)$ is polyhedral.
%One can view $f_{\alpha}$ as a function of $\alpha$ defined on $\BR(\L_n^+)$. It is easy to see that $f_{\alpha+\beta}\leq f_{\alpha}f_{\beta}$ and $f_{r\alpha}=(f_{\alpha})^{r}$ for any $r>0$. We recall that $\BR(\L_2^+)$ are exactly the nonnegative multiples of 
%\[\alpha=(\alpha_{11},\alpha_{12},\alpha_{22})=(1,-2,1)\]
%with the optimal constant $C_\alpha=1$, thus $f_{\alpha}=1$ is a constant function on $\BR(\L_2^+)$. In the next smallest case, the function $f_{\alpha},~~~\alpha \in \BR(\L_3^+)$ is quite interesting. 

%Firstly, we state a supplementary lemma. Recall that there is a scaling action $(\RR_{>0})^n$ on $\L_n^+$, where $(c_1,\dots, c_n)\cdot (p_{ij})_{ij} = (c_ic_jp_{ij})_{ij}$.
%Recall also that $\pi:\RR^{\{(i,j):1\le i\le j\le n\}}\to \RR^{\binom{[n]}{2}}$ is a projection which restricts to a bijection between $\BR(\L_n^+)$ and $\BR'(\underline{\L_n^+})$. %Similarly, for $\alpha \in \BR'(\underline{\L}_n^+)$, we define an optimal upper bound
%\[f^*_{\alpha}\coloneqq \inf \left\{C_\alpha>0 \text{ such that }\forall p\in \underline{\L}_n^+, \prod_{1\le i< j\le n}p_{ij}^{\alpha_{ij}}\le C_\alpha\right\}.\]

%\begin{lemma}
%\label{lem:same-constant}
%Given any $\alpha \in \BR(\L_n^+)$, the optimal constants satisfy $f_{\alpha}=f_{\pi(\alpha)}$.
%\end{lemma}
%\begin{proof}
%Since normalized Lorentzian matrices (Lorentzian matrices with all-ones diagonal) form a subset of Lorentzian matrices, we have  $f_{\alpha} \geq f_{\pi(\alpha)}$. 

%For any $\alpha \in \BR(\L_n^+)$ the following balancing condition holds
%\[
%	\text{for all } i\in [n],\; 2\alpha_{ii}+\sum_{j\neq i}\alpha_{ij}=0.
%\]
%Indeed, otherwise the bounded ratio $\alpha$ attains arbitrary large values on the orbit of the scaling action $(1,\dots, t,\dots, 1)\cdot (p_{ij})_{ij}$, which contradicts to $\alpha \in \BR(\L_n^+)$. The balancing conditions imply that the value of a bounded ratio $\alpha$ is constant on a torus orbit. Clearly, any Lorentzian matrix is in the orbit of some normalized Lorentzian matrix. The claim then follows. 
%\end{proof}

%Due to the above lemma we study optimal constants for reduced bounded ratios. By Theorem~\ref{thm:BR}, the cone $\BR'(\underline{\L_3^+})$ has primitive ratios 
%\begin{equation}\label{eq:r123}
 %   r_1=(-1,-1,1), \quad r_2=(-1,1,-1),\quad \text{and}\quad r_3=(1,-1,-1).
%\end{equation}
%$f^*_{\alpha}, ~~~\alpha \in \BR'(\underline{\L}_3^+)$ instead of $f_{\alpha}, ~~~\alpha \in \BR(\L_3^+)$. We will abuse notation, and write $f_{\alpha}$ instead of $f^*_{\alpha}$ from now on.

%The cone $\BR'(\underline{\L}_3^+)$ has three extremal rays corresponding to the ratios 
%\[\frac{p_{23}}{p_{12}p_{13}},~~~\frac{p_{13}}{p_{12}p_{23}},~~~\frac{p_{12}}{p_{13}p_{23}}\]
%which are projections of the triangular inequalities under the map $\pi$. Having fixed lexicographical order on the coordinates the three corresponding vectors are
%\[r_1=(-1,-1,1), \quad r_2=(-1,1,-1), \quad r_3=(1,-1,-1).\]
%
%Since $f_{r\alpha}=f^{r}_{\alpha}$ for any $r>0$, it suffices to determine values of $f_{\alpha}$ on the subset of $\BR'(\underline{\L_3^+})$ given by nonnegative linear combinations 
%\[\alpha=ar_1+br_2+cr_3, \quad \text{where} \quad a+b+c=1, \quad a,b,c \geq 0.\]
%

%We now prove Theorem~\ref{mth:entropy}.

\begin{proof}[Proof of Theorem~\ref{mth:entropy}]
We consider the triangular section of the cone of reduced bounded ratios 
\[
\left\{a \cdot \alpha^{23|1}+b \cdot \alpha^{13|2}+c \cdot \alpha^{12|3}  \;\middle|\;  \text{$a,b,c$ are nonnegative numbers that sum to $1$}\right\} \subseteq \underline{\BR}(\L_3^+).
\]
The condition $a^2+b^2+c^2-2ab-2ac-2bc=0$ defines an inscribed circle in the given triangle, where the three tangent points are given by
\[
(a,b,c) \ = \ (1/2,1/2,0), \ \ (1/2,0,1/2), \ \ (0,1/2,1/2).
\]
These points are the midpoints of the sides of the triangle, corresponding to the Alexandrov--Fenchel inequalities. 

The inscribed circle splits the triangle in four regions. The part within the circle is defined by 
\[
a^2+b^2+c^2-2ab-2ac-2bc\leq 0.
\]
Outside the circle, the connected component containing the vertex $\alpha^{23|1}$ is defined by
\[
a^2+b^2+c^2-2ab-2ac-2bc > 0,  \ \ a \ge b, \ \ a \ge c.
\]
The other two regions outside of the inscribed circle are described similarly. 
%\begin{enumerate}
%    \item $a^2+b^2+c^2-2ab-2ac-2bc >0, \quad a>b, a>c$, which is connected part of the triangle outside the inscribed circle and containing vertex $a=1,~b=c=0$.
%    \item $a^2+b^2+c^2-2ab-2ac-2bc >0, \quad b>a, b>c$, which is connected part of the triangle outside the inscribed circle and containing vertex $b=1,~a=c=0$.
%    \item $a^2+b^2+c^2-2ab-2ac-2bc >0, \quad c>a, c>b$, which is connected part of the triangle outside the inscribed circle and containing vertex $c=1,~a=b=0$.
%    \item $a^2+b^2+c^2-2ab-2ac-2bc\leq 0$, which is the part within the circle and includes its boundary. 
%\end{enumerate}

We may reformulate the objective as the problem of   finding the global maximum of 
\[
R(p_{12},p_{13},p_{23})\coloneqq(p_{12})^{c-b-a}(p_{13})^{b-c-a}(p_{23})^{a-b-c},
\]
subject to certain constrains on $p_{12},p_{13},p_{23}$. More precisely, we want to find the global maximum of $R(p_{12},p_{13},p_{23})$ in terms of the parameters $a,b,c$  when $p_{12},p_{13},p_{23}\geq 1$ and
\[
\det\begin{pmatrix}
	1 & p_{12} & p_{13} \\
	p_{12} & 1      & p_{23} \\
	p_{13} & p_{23} & 1
\end{pmatrix}= 
1+2p_{12}p_{13}p_{23}-p^2_{12}-p^2_{13}-p^2_{23}\geq 0.
\]
These inequalities are equivalent to the condition that  $3 \times 3$ matrix $(p_{ij})$ being Lorentzian.

We apply Lagrange's method to $R(p_{12},p_{13},p_{23})$ separately on the three regions outside the inscribed circle. On the region containing $\alpha^{23|1}$,  one finds that the interior of the domain has no critical points, and thus the maximum is attained on the boundary. 
The interesting case is when 
\[
1+2p_{12}p_{13}p_{23}-p^2_{12}-p^2_{13}-p^2_{23}=0.
\]
Since $p_{12},p_{13},p_{23} \ge 1$, the condition is equivalent to 
\[
p_{23}=p_{12}p_{13}+\sqrt{(p_{12}-1)(p_{13}-1)}.
\]
Thus, setting $x=p^{-2}_{12}$ and $y=p^{-2}_{13}$, the question reduces to finding the global maximum of 
%\begin{equation}\label{eq:xy function}
\[
r(x,y) \coloneq x^{b}y^{c}\left(1+\sqrt{(1-x)(1-y)}\right)^{a-b-c} \ \ \text{when $0\leq x,y \leq 1$}.
\]
%\end{equation}
The unique critical point of $r(x,y)$ has coordinates 
\[
x_0=4ab/(a+b-c)^2  \ \ \text{and} \ \  y_0=4ac/(a-b+c)^2,
\]
and they satisfy $0 \le x_0, y_0 \le 1$ because $a^2+b^2+c^2-2ab-2ac-2bc >0$. Evaluating, we get
\[
r(x_0,y_0)=2\cdot a^a\cdot b^b \cdot c^c \cdot (2a-1)^{2a-1} \cdot (1-2b)^{2b-1}\cdot (1-2c)^{2c-1}.
\]
%, which are in the domain for the given restrictions on $a,~b,~c.$ %Then, 
%\[f_{\alpha}=\frac{2a^ab^bc^c(2a-1)^{2a-1}}{(1-2c)^{1-2c}(1-2b)^{1-2b}}, \quad \text{for} \quad a^2+b^2+c^2-2ab-2ac-2bc >0, \quad a>b, a>c.\]
When $a^2+b^2+c^2-2ab-2ac-2bc=0$, we have $(a+b-c)^2=4ab$ and $(a-b+c)^2=4ac$. In this case, the critical point is $(1,1)$ and $r(x,y)$ has value $1$ at the critical point. 
%\[(a+b-c)^2=4ab, \quad (a-b+c)^2=4ac\quad \text{and}\quad (a-b-c)^2=4bc. 
%\]
Thus, when $(a,b,c)$ lies in the the boundary of the inscribed circle, $f(a,b,c)=1$. By the log-convexity of $f(a,b,c)$, 
 we conclude that $f(a,b,c)\leq 1$ inside the circle. On the other hand, the matrix with all entries $1$ is Lorentzian, and any ratio for this matrix equals $1$. Thus, $f(a,b,c)=1$ inside the circle.
\end{proof}


We compare the optimal bounding constants $f(a,b,c)$ for $\L_3^+$ with the corresponding constants $f_p(a,b,c)$ for $\Delta_3^+(\T_p)$, which is piecewise log-linear. 
%We write $g_\alpha^p$ for the optimal constants $f_\alpha^{\Delta(\T_p)_n^+}$.
%Again, we work with reduced bounded ratios. 
%The cone $\underline{\BR}(\Delta_n^+(\T_p))$ has the same three extremal rays generated by $r_1,~r_2,~r_3$ as in \eqref{eq:r123}.
\begin{proposition}
Consider the reduced bounded ratio 
$\alpha=a \cdot \alpha^{23|1}+b \cdot \alpha^{13|2}+c \cdot \alpha^{12|3}$ on $\Delta_n^+(\T_p)$, where $a,b,c$ are nonnegative numbers. 
The optimal bounding constant for $\alpha$ is given by
\[
f_p(a,b,c)=\begin{cases}
2^{ap} & \text{if} \; a>b+c,\\
2^{bp} & \text{if} \; b>a+c,\\
2^{cp} & \text{if} \; c>a+b,\\
2^{\frac{p(a+b+c)}{2}} & \text{if otherwise}.
\end{cases}
\]	    
\end{proposition}
\begin{proof}
Let $q_{ij}=\log_2(p_{ij}),$ where $i,j$ are any two distinct indices from $\{1,2,3\}$. In log coordinates, the problem reduces to the maximization of the linear function
\[
\ell(a,b,c)\coloneq a(q_{23}-q_{12}-q_{13})+b(q_{13}-q_{12}-q_{23})+c(q_{12}-q_{13}-q_{23}), 
\]
where the domain given by the linear inequalities 
\begin{align*}
q_{23} &\leq p + q_{12} + q_{13}, & \
q_{13} &\leq p + q_{12} + q_{23}, & \
q_{12} &\leq p + q_{13} + q_{23}, \\
0 &\leq p + q_{12}+q_{12}, & \
0 &\leq p + q_{13}+q_{13}, & \
0 &\leq p + q_{23}+q_{23}.
\end{align*}
The maximum of $\ell(a,b,c)$ is attained at one of the vertices of the domain. The solution to the system contains four vertices
\[
 (0,-p/2,-p/2), \ \ (-p/2,0,-p/2), \ \ (-p/2,-p/2,0), \ \ \text{and} \ \ (-p/2,-p/2,-p/2),
\]
corresponding to the four cases in the displayed equation for $f_p(a,b,c)$.
%Then
%\[
%r(a,b,c)=\begin{cases}
%ap, &\text{if} \quad a>b+c,\\
%bp, & \text{if} \quad b>a+c,\\
%cp, & \text{if} \quad c>a+b,\\
%\frac{p(a+b+c)}{2} & \text{if} \; a,b,c \text{ form a possibly degenerate triangle}.
%\end{cases}
%\]	  
%Then $g_{\alpha}^p=2^{r(a,b,c)}$, which completes the proof.
\end{proof}



%%For any $\mathbb{T}_2$-matrix $(p_{ij})$, we have
%\[
%\frac{p_{12}p_{13}p_{23}p_{45}p_{44}p_{55}}{p_{14}p_{15}p_{24}p_{25}p_{34}p_{35}} \le 64.
%\]
%The equality is achieved by the $\mathbb{T}_2$-matrix
%\[
%\begin{pmatrix}
%0 & 2 & 2 & 1 & 1 \\
%2 & 0 & 2 & 1 & 1 \\
%2 & 2 & 0 & 1 & 1 \\
%1 & 1 & 1 & 2 & 2 \\
%1 & 1 & 1 & 2 & 2
%\end{pmatrix}.
%\]
%(State this for general $p$, including $0$)
%We have
%\[
%\sqrt{p_{44}p_{13}} \le 2 \sqrt{p_{14}p_{34}}, \quad \sqrt{p_{55}p_{23}} \le 2 \sqrt{p_{25}p_{35}}, \ \ \text{and hence} \ \ \sqrt{p_{44}p_{55}p_{13}p_{23}} \le 4 \sqrt{p_{14}p_{34}p_{25}p_{35}}.
%\]
%Similarly, we have
%\[
%\sqrt{p_{44}p_{23}} \le 2 \sqrt{p_{24}p_{34}}, \quad \sqrt{p_{55}p_{13}} \le 2 \sqrt{p_{15}p_{35}}, \ \ \text{and hence} \ \ \sqrt{p_{44}p_{55}p_{13}p_{23}} \le 4 \sqrt{p_{24}p_{34}p_{15}p_{35}}.
%\]
%Now consider the triangle inequality
%\[
%\sqrt{p_{12}p_{45}} \le \sqrt{p_{14}p_{25}}+\sqrt{p_{15}p_{24}}.
%\]
%This gives
%\[
%\sqrt{p_{12}p_{13}p_{23}p_{44}p_{55}p_{45}} \le \sqrt{p_{44}p_{55}p_{13}p_{23}p_{14}p_{25}}+\sqrt{p_{44}p_{55}p_{13}p_{23}p_{15}p_{24}} \le 8 \sqrt{p_{14}p_{15}p_{24}p_{25}p_{34}p_{35}}.
%\]
%\end{remark}

\section{Rank 2 Lorentzian matrices}\label{sec:rank2}
%\textcolor{blue}{Here we explain that}
%\[
%\text{rank $2$ Lorentzian quadratic forms}=\text{product of two nonnegative linear forms}.
%\]
%
%\textcolor{blue}{Theorem: The cone of bounded ratios of these is also the dual of the cut cone.}
%
%\textcolor{blue}{The corresponding conjecture can be strengthened as follows.}
%
%\textcolor{blue}{Conjecture: The bounded ratios for rank $2$ Lorentzian quadratic forms give subtraction-free expressions.}
%
%\textcolor{blue}{For example, (show a simple example.)}
%The space of Lorentzian matrices is subtle to describe. 

We show that the cone of bounded ratios on Lorentzian matrices is equal to the cone of bounded ratios on rank $2$ Lorentzian matrices. 
The space of rank $2$ Lorentzian matrices admits a natural parametrization, which allows us 
 to formulate a stronger notion of positivity called the \emph{subtraction-freeness}. Throughout, we work with $n \times n$ symmetric matrix $M$ with real entries. %In this section, we study the set of bounded ratios for rank-2 Lorentzian matrices and propose several conjectures.

\begin{lemma}\label{lem:quadr1} 
The quadratic form $Q(x)=x^{T}Mx$ factors into a product of two linear forms if and only if there is a rank at most $1$ matrix $B$  such that $M=\frac{1}{2}(B+B^{T})$.    
\end{lemma}
\begin{proof}
If the quadratic form $Q(x)$ factors into linear forms, we can rewrite
\[
Q(x)=x^{T}Mx=(x^{T}a)(b^{T}x)=x^{T}(ab^{T})x=x^T(ba^T)x,
\]
where $a$ and $b$ are column vectors. Then  $B=ab^{T}$ is a rank at most $1$ matrix and $Q(x)=x^{T}\left(\frac{B+B^{T}}{2}\right)x$.
Conversely, if $M=\frac{1}{2}(B+B^{T})$ for a rank at most $1$ matrix $B$, then 
\[Q(x)=x^{T}Mx=x^{T}\left(\frac{B+B^{T}}{2}\right)x=x^{T}Bx.\]
Since $B$ has rank at most $1$, there are real vectors $a$ and $b$ such that
$x^{T}Bx=x^{T}(ab^{T})x=(x^{T}a)(b^{T}x)$. 
\end{proof}

\begin{lemma}\label{lem:rank2} 
 There is a rank at most $1$ matrix $B$ such that $M=\frac{1}{2}(B+B^{T})$ if and only if $M$ has rank at most $2$ and all $2\times2$ principal minors of $M$ are nonpositive.
\end{lemma}
\begin{proof}
If  $M=\frac{1}{2}(B+B^{T})$ for a rank at most $1$ matrix $B$, then, by Lemma \ref{lem:quadr1}, there is a factorization 
\[
x^{T}Mx=\left(\sum_{1\leq i\leq n}a_ix_i\right)\left(\sum_{1\leq i\leq n}b_ix_i\right).
\]
Thus, $M$ has rank at most $2$ and any $2 \times 2$ principal minor of $M$ is nonpositive:
\[(2a_ib_i)(2a_jb_j)-(a_ib_j+a_jb_i)^2=-a_i^2b_j^2+2a_ia_jb_ib_j-a_j^2b_i^2 \leq 0.\]
%One can also verify directly that any minor of size $3$ is identically zero. %Indeed, for any two triples $i,j,k$ and $l,m,n$, possibly overlapping, the following identity holds
%\[
%\det
%\begin{pmatrix}
%a_{i}b_{l}+a_{l}b_{i} & a_{i}b_{m}+a_{m}b_{i} & a_{i}b_{n}+a_{n}b_{i} \\
%a_{j}b_{l}+a_{l}b_{j} & a_{j}b_{m}+a_{m}b_{j} & a_{j}b_{n}+a_{n}b_{j} \\
%a_{k}b_{l}+a_{l}b_{k} & a_{k}b_{m}+a_{m}b_{k} & a_{k}b_{n}+a_{n}b_{k} \\
%\end{pmatrix}\equiv 0
%\]
%Thus, $M$ is of rank at most $2$.

For the converse, suppose that $M$ has rank at most $2$ and that all $2\times2$ principal minors of $M$ are nonpositive. We want to find a matrix $B=ab^T$ 
for some real vectors $a,b$ such that $M=\frac{1}{2}(B+B^{T})$. If $M$ is rank at most 1, $M=uu^T$ for some $u$ and we take $a=b=u$. Now suppose $M$ is rank 2. 
Since $M$ has all $2\times 2$ principal minors nonpositive, $M$ cannot be positive or negative semidefinite, so $M$ must have eigenvalues $\lambda_1>0>\lambda_2$. 
Suppose $u$ and $v$ are orthonormal eigenvectors with these eigenvalues. Then $M=\lambda_1 uu^T+\lambda_2vv^T$. 
It then follows that we can take $a=(\lambda_1)^{1/2}u+(-\lambda_2)^{1/2}v$ and $b= (\lambda_1)^{1/2}u-(-\lambda_2)^{1/2}v$.
\end{proof}


\begin{proposition}
For any nonnegative vector $(a_i)_{1 \le i \le n}, (b_i)_{1 \le i \le n}$, the Hessian of the quadratic form $\left(\sum_{1\leq i\leq n}a_ix_i\right)\left(\sum_{1\leq i\leq n}b_ix_i\right)$ is Lorentzian, and every Lorentzian matrix of rank at most $2$ is of this form.
\end{proposition}

\begin{proof}
Let $M$ be the Hessian of the quadratic form  $Q=\left(\sum_{1\leq i\leq n}a_ix_i\right)\left(\sum_{1\leq i\leq n}b_ix_i\right)$, where $a,b$ are nonnegative real vectors. 
By Lemma \ref{lem:quadr1} and Lemma \ref{lem:rank2}, we know that $M$ is of rank at most $2$ and all $2\times2$ principal minors of $M$ are nonpositive. Since $a,b$ are nonnegative, the entries of $M$ are nonnegative. Thus, the symmetric matrix $M$ is a Lorentzian matrix of rank at most $2$.

Conversely, for any nonzero Lorentzian matrix $M$ of rank at most $2$, by Lemma \ref{lem:quadr1} and Lemma \ref{lem:rank2}, the corresponding quadratic form $x^{T}Mx$ factors into a product 
\[x^{T}Mx=\left(\sum_{1\leq i\leq n}a_ix_i\right)\left(\sum_{1\leq i\leq n}b_ix_i\right) \ \ \text{for some real vectors $a,b$.}\]  
It remains to show that vectors $a,b$ can be chosen nonnegative. Since the entries of $M$ are nonnegative, $a_ib_i\geq 0$ for all $1\leq i\leq n$. %Assume that $M$ is nonzero. %Replacing all $a_i$ and $b_i$ by its negative if necessary, and after an appropriate permuting the indices, we may assume that $a_1,b_1\geq 0$ and $(a_i, b_i)\neq (0,0)$. %,\ otherwise we can switch sign of all coefficients in both linear forms. 

First, consider the case when $a_ib_i=0$ for $1\leq i\leq n$. Since $M$ is nonzero, replacing all $a_i$ and $b_i$ by its negative and permuting the indices if necessary, we may suppose that $a_1>0$, which in turn implies that $b_1=0$. For any $b_i\neq 0$, since $a_1b_i+a_ib_1\geq0$, we have $b_i>0$. Thus, all nonzero $b_i$ are positive. Since $M$ is nonzero, one of the $b_i$ must be positive. Then we apply the same argument to conclude that all nonzero $a_i$ are positive. 

When $a_ib_i> 0$ for some $i$, without loss of generality, we may assume that $a_1>0$ and $b_1>0$. Since $a_1b_i+a_ib_1$ is nonnegative as well as $a_ib_i\geq0$ for any $i$, we conclude that all $a_i$ and $b_i$ are nonnegative. 
\end{proof}

Let $\L_n^+(2)$ denote the set of $n\times n$ Lorentzian matrices of rank at most $2$.
Since $\L_n^+(2)\subseteq \L_n^+$, we have 
$\underline{\BR}\bigl(\L_n^+\bigr) \subseteq \underline{\BR}\bigl(\L_n^+(2)\bigr)$.
We show that these cones coincide and they are dual to the cut cone $\Cut_n$.

\begin{proposition}\label{thm:BR_rk2}
We have 
	$
    \underline{\BR}\bigl(\L_n^+\bigr)=
    \underline{\BR}\bigl(\L_n^+(2)\bigr)$.
    %\big\{\alpha\in \RR^{\binom{n}{2}}: \alpha\cdot \beta \le 0 \ \text{ for all } \beta\in \Cut_n\big\}$.
\end{proposition}

\begin{proof}
Theorem~\ref{mth:BR} says that $\Cut_n$ is dual to  $\underline{\BR}\bigl(\L_n^+\bigr)$. Thus, the dual of  $\Cut_n$ is contained in $\underline{\BR}\bigl(\L_n^+(2)\bigr)$. 
We show the opposite inclusion.





%Let $E=\{\exp(c\delta(S)):c\ge 0,\emptyset\neq S\subseteq[n]\}$. 
Let  $E$ be the set of entry-wise exponentials of extremal rays of $\Cut_n$, viewed as $n \times n$ symmetric matrices with diagonal entries $1$.
%By Lemma \ref{lem:convhullT0} we know that the convex hull of $\log \underline{\Delta(\T_0)_n^+}$ is  $\Cut_n$. 
Since every element of $E$ is the tree metric of a star tree, it follows that $E \subseteq \underline{\Delta}_n^+(\T_0) \subseteq \underline{\L}_n^+$. Since every matrix  in $E$ has rank at most 2, we have
$E \subseteq \underline{\L}_n^+(2)$.
Thus, any reduced bounded ratio $\alpha$ on $\L_n^+(2)$ satisfies $\alpha \cdot \beta \le 0$ for all $\beta \in \log E$, and hence for all $\beta \in \Cut_n$. 
 %$\BR'\bigl(\underline{\L_n^+|_{\mathrm{rk} 2}}\bigr) \subseteq \BR'(E)$.
%It remains to show that $\BR(E)$ is equal to the dual of  $\Cut_n$. In fact, this equality follows from the fact that the convex hull of $E$ is $\Cut_n$ (see the proof of Lemma~\ref{lem:convhullT0} (iii)).
%\begin{align*}
%\BR(E)
%&=\left\{ \alpha\in \RR^{\binom{[n]}{2}}:\exists c> 0\text{ such that for all } S, \sum_{1\le i< j\le n}\alpha_{ij}\delta(S)_{ij}\le 	\log c\right\}\\
%&=\left\{ \alpha\in \RR^{\binom{[n]}{2}}:\exists c> 0\text{ such that }\forall q\in \Cut_n, \sum_{1\le i< j\le n}\alpha_{ij}q_{ij}\le 	\log c\right\}\\
%&= \left\{ \alpha\in \RR^{\binom{[n]}{2}}:\forall q\in \Cut_n, \sum_{1\le i< j\le n}\alpha_{ij}q_{ij}\le 	0\right\}.
%\end{align*}
%The first equality follows from the definition of $\BR$, the second follows from the fact that $\Cut_n$ has extremal rays given by extremal tree metrics $\delta(S)$, and the last is because $\Cut_n$ is a finitely generated convex polyhedral cone. This concludes the proof. 
\end{proof}

We close with a statement that strengthens Conjecture~\ref{conj:intro-at-most-2} for rank $2$ Lorentzian matrices.

\begin{conjecture}\label{conj:sub-free}
Let $\alpha \in \BR(\L_n^+)$ be any integral bounded ratio and $(p_{ij})_{1\le i\le j\le n}$ be the Hessian matrix of the quadratic form $\left(\sum_{1\leq i\leq n}a_ix_i\right)\left(\sum_{1\leq i\leq n}b_ix_i\right)$. Then %the ratio $\prod_{1\le i\leq j\le n}p_{ij}^{\alpha_{ij}}$ is \emph{subtraction-free}, i.e.   
\[
2^{\sum_{i=1}^n\alpha_{ii}}\Biggl(\prod_{\alpha_{ij}<0}p_{ij}^{-\alpha_{ij}}\Biggr)-\Biggl(\prod_{\alpha_{ij}>0}p_{ij}^{\alpha_{ij}}\Biggr)\]
	is a polynomial with nonnegative coefficients in $a_i$ and $b_i$. 
\end{conjecture}

%\begin{example}
For example, for $(a_1x_1+a_2x_2+a_3x_3)(b_1x_1+b_2x_2+b_3x_3)$ with $a_i,b_i \ge 0$, we have 
\[
\frac{p_{11}p_{23}}{p_{12}p_{13}}=\frac{(2a_1b_1)(a_2b_3+a_3b_2)}{(a_1b_2+a_2b_1)(a_1b_3+a_3b_1)} \le 2,
\]
and, in fact, the difference $2p_{12}p_{13}-p_{11}p_{23}=2a^2_1b_2b_3+2a_2a_3b_1^2$
has nonnegative coefficients.
%\end{example}
Computations show that, for $n \le 6$, every integral bounded ratio on $\L_n^+$ is a nonnegative integral linear combination of primitive bounded ratios on $\L_n^+$, and, every primitive bounded ratios on on $\L_n^+$ satisfy Conjecture~\ref{conj:sub-free}. Since the validity of the conjecture is preserved under taking nonnegative integral linear combinations, this verifies Conjecture~\ref{conj:sub-free} for $n \le 6$.
%Conjecture \ref{conj:sub-free} holds for all primitive bounded ratios up to $n\leq 7$. 
%\begin{remark}
For a similar subtraction-free phenomenon for bounded ratios of Pl\"ucker coordinates over totally positive Grassmannians was observed in \cite[Conjecture 37]{BoFr}. For a connection to cluster algebras,  see \cite{Gekhtman-Greenberg-Soskin}.   
%\end{remark}
%Before we suggest a conjecture on upper bound for bounded ratios over set of all Lorentzian matrices we state the following lemma. 
%\begin{lemma}
%\label{lem:alpha-negative}
%Suppose  $0\neq \alpha \in \BR\bigl(\L_n^+\bigr)$. Then $\sum_{i<j}\alpha_{ij}<0$.    
%\end{lemma}
%\begin{proof}
%Given $p \in \L_n^+$, one can fix all entries of $p$, except one diagonal entry $p_{ii}$ for some $1 \leq i \leq n$. When the entry $p_{ii}$ is replaced by a sufficiently small and positive number, the matrix remains in $\L_n^+$. Thus, $\alpha$ is a bounded ratio implies that $\alpha_{ii}\geq 0$ for all $1 \leq i \leq n$. Since $\sum_{i\leq j}\alpha_{ij}=0$ it remains to prove that at least one $\alpha_{ii}>0$.

%Suppose  that for all $i$, $\alpha_{ii}=0$ and thus $\sum_{\{j:j\neq i\}}\alpha_{ij}=0$. Then $\sum_{i\neq j}\alpha_{ij}=0$. We show that this implies $\alpha$ is the zero vector. Because otherwise there exists $k\neq l$ such that $\alpha_{kl}<0$. Let $S=\{k,l\}$. Then
%\begin{align*}
%	\alpha\cdot \delta(S) 
%	=&\sum_{i\neq j} \alpha_{ij}\delta(S)_{ij}\\
%	=&\sum_{\substack{i\neq j\\\{k,l\}\cap\{i,j\}=k}}\alpha_{ij} + \sum_{\substack{i\neq j\\\{k,l\}\cap\{i,j\}=l}}\alpha_{ij} \\
%	=&\sum_{\{j:j\neq k\}}\alpha_{kj} + \sum_{\{i:i\neq l\}}\alpha_{il}-2\alpha_{kl},\\
 %       =&-2\alpha_{k,l}>0
%\end{align*}
%which, by Theorem~\ref{mth:BR}, implies that $\alpha\notin \BR(\L_n^+)$, a contradiction.
%\end{proof}
%Since $f_{r\alpha}=(f_{\alpha})^{r}$ for any $\alpha \in \BR(\L_n^+)$, it suffices to consider normalized 
%$\alpha$, namely those with $\sum_{i<j}\alpha_{ij}=-1$.
%it is natural to consider only normalized $\alpha$, i.e. such $\alpha$ that $\sum_{i<j}\alpha_{ij}=-1$.
%Conjecture \ref{conj:sub-free} implies that for a normalized bounded ratio $\alpha$ 
%$$\prod_{1\le i\leq j\le n}p_{ij}^{\alpha_{ij}}\leq2$$
%on $\L_n^+|_{\mathrm{rk} 2}$. In general, switching to a larger set $\L_n^+$ may increase an upper bound for a given ratio. 
%We conjecture that 2 remains an upper bound even when $p\in \L_n^+$.

%\begin{conjecture}\label{conj:at-most-2}
%Consider any normalized $\alpha \in \BR(\L_n^+)$. Then for any $p\in\L_n^+$,
%$$\prod_{1\le i\leq j\le n}p_{ij}^{\alpha_{ij}}\leq2.$$   %\end{conjecture}
%By Theorem~\ref{mth:pentagonal}, we know that Conjecture \ref{conj:at-most-2} holds for $n\leq 5$.







%The determinant of $M$ is identically zero, and all the $4 \times 4$ principal minors of $M$ are identically zero. The ten $3 \times 3$ principal minors of $M$ are
%\begin{align*}
%\det M_{123} &= 2,      & \det M_{124} &= 0,   & \det M_{125} &= 8,      & \det M_{134} &= 2x^2,   & \det M_{135} &= 2x^2,\\
%\det M_{145} &= 8x^2,   & \det M_{234} &= 8,   & \det M_{235} &= 0,      & \det M_{245} &= 32,     & \det M_{345} &= 8x^2.
%\end{align*}
%and the ten $2 \times 2$ principal minors of $M$ are
%\begin{align*}
%\Delta_{12}&=-1, & \Delta_{13}&=-1, & \Delta_{14}&=-x^2, & \Delta_{15}&=-(2+x)^2, & \Delta_{23}&=-1,\\
%\Delta_{24}&=-4, & \Delta_{25}&=-4, & \Delta_{34}&=-(2+x)^2, & \Delta_{35}&=-x^2, & \Delta_{45}&=-16\,(1+2x),
%\end{align*}

\bibliography{biblio}
%\bibliographystyle{plain}
\bibliographystyle{amsalpha}

\end{document}
